\chapter*{Exercises}
\phantomsection
\addcontentsline{toc}{chapter}{Exercises}
\setcounter{exercisegroup}{1}
\edef\CurrentExerciseGroup{\arabic{exercisegroup}}
\section*{\texorpdfstring{\S\CurrentExerciseGroup}{Section \CurrentExerciseGroup}}
\phantomsection
\addcontentsline{toc}{section}{\texorpdfstring{Exercises for \S\CurrentExerciseGroup}{Exercises for Section \CurrentExerciseGroup}}

\begin{enumerate}
\def\labelenumi{\arabic{enumi})}
\item
  Consider the locally compact spaces T, T' identical to the interval {[}0, 1{]}, and the measures \(\mu, \mu'\) on T and T', respectively, identical to the Lebesgue measure. Let F be a Hilbert space admitting a countable orthonormal basis, arranged in a double sequence \((\mathbf{e}_{mn})\) ; let F' = F be its dual. Set \(\mathbf{u}_m(t) = \mathbf{e}_{mn}\) for \((n-1)2^{-m} \leqslant t < n2^{-m}\) and \(1 \leqslant n \leqslant 2^m\) ; set \(\mathbf{f}(t, t') = 2^m \mathbf{u}_m(t)\) for \(2^{-m} < t' \leqslant 2^{-m+1}\) , and finally \(\mathbf{f}(t, 0) = 0\) and \(\mathbf{f}(1, t') = 0\) . Show that \(\mathbf{f}\) is scalarly integrable for the product measure \(\mu \otimes \mu'\) , but that the function \(t' \mapsto \mathbf{f}(t, t')\) is not scalarly \(\mu'\) -integrable for any \(t \in T\) .
\item
  Let \(\mathrm{T}, \mathrm{X}, \mathrm{Y}\) be three locally compact spaces, with \(\mathrm{Y}\) having a countable base. Let \(\mu\) be a measure \(\geqslant 0\) on \(\mathrm{T}\) , \(t \mapsto \lambda_t (t \in \mathrm{T})\) a \(\mu\) -adequate \(^1\) family of measures \(\geqslant 0\) on \(\mathrm{X}\) , and set \(\nu = \int \lambda_t d\mu(t)\) . Let \(x \mapsto \rho_x (x \in \mathrm{X})\) be a \(\nu\) -adequate family of measures \(\geqslant 0\) on \(\mathrm{Y}\) . Suppose, moreover, that one of the following two conditions is satisfied: a) \(\mathrm{X}\) is countable at infinity;
\end{enumerate}

\begin{enumerate}
\def\labelenumi{\alph{enumi})}
\setcounter{enumi}{1}
\tightlist
\item
  \(\nu\) is bounded.
\end{enumerate}

Under these conditions, show that there exists a locally \(\mu\) -negligible set \(N'\) such that, for every \(t \notin N'\) , the family \(x \mapsto \rho_x\) is \(\lambda_t\) -adequate; moreover, the family \(t \mapsto \int \rho_x d\lambda_t(x)\) , defined for \(t \notin N'\) , is \(\mu\) -adequate, and

\[
\int d \mu (t) \int \rho_ {x} d \lambda_ {t} (x) = \int \rho_ {x} d \nu (x).
\]

(To prove that \(x \mapsto \rho_x\) is scalarly \(\lambda_t\) -integrable locally almost everywhere, make use of Lemma 1 of §3, No.~1, using, in case \(b\) ), the fact that \(\lambda_t\) is bounded locally almost everywhere; same method for seeing that \(t \mapsto \int \rho_x d\lambda_t(x)\) is scalarly \(\mu\) -integrable; use also Prop. 13 of No.~5.)

\begin{enumerate}
\def\labelenumi{\arabic{enumi})}
\setcounter{enumi}{2}
\tightlist
\item
  Let S, T, Y be three locally compact spaces, countable at infinity, with Y having a countable base. Let \(\rho\) (resp. \(\sigma\) ) be a positive measure on S (resp. T), \(\nu = \rho \otimes \sigma\) , and let \((\lambda_{s,t})_{(s,t) \in \mathbb{S} \times \mathbb{T}}\) be a \(\nu\) -adequate \(^{1}\) family of positive measures on Y. Then, there exists a \(\rho\) -negligible set N such that, for every \(s \notin N\) , the family \((\lambda_{s,t})_{t \in \mathbb{T}}\) is \(\sigma\) -adequate; the family \(s \mapsto \int \lambda_{s,t} d\sigma(t)\) (defined almost everywhere) is \(\rho\) -adequate, and
\end{enumerate}

\[
\iint \lambda_ {s, t} d \rho (s) d \sigma (t) = \int d \rho (s) \int \lambda_ {s, t} d \sigma (t).
\]

(Apply Exer. 2 to the space \(\mathrm{X} = \mathrm{S} \times \mathrm{T}\) , on noting that \(\nu = \int (\varepsilon_s \otimes \sigma) d\rho(s)\) .)

\begin{enumerate}
\def\labelenumi{\arabic{enumi})}
\setcounter{enumi}{3}
\tightlist
\item
  Let \(\mathrm{T}, \mathrm{X}, \mathrm{Y}\) be three locally compact spaces, with \(\mathrm{X}\) being countable at infinity. Let \(\mu\) be a positive measure on \(\mathrm{T}\) , \((\lambda_t)_{t \in \mathrm{T}}\) a \(\mu\) -adequate family of positive measures on \(\mathrm{X}\) , and \(\nu = \int \lambda_t d\mu(t)\) . Assume that one of the following two conditions holds:
\end{enumerate}

\begin{enumerate}
\def\labelenumi{\alph{enumi})}
\item
  Y has a countable base;
\item
  the measure \(\nu\) is bounded.
\end{enumerate}

Under these conditions, if \(\pi\) is a \(\nu\) -proper mapping of X into Y, then the set of \(t \in \mathbf{T}\) such that \(\pi\) is not \(\lambda_t\) -proper is locally \(\mu\) -negligible; moreover, the family \((\pi(\lambda_t))_{t \in \mathbf{T}}\) of positive measures on Y (defined locally almost everywhere in T) is \(\mu\) -adequate, and

\[
\pi \left(\int \lambda_ {t} d \mu (t)\right) = \int \pi (\lambda_ {t}) d \mu (t).
\]

(In case \(a\) ), apply Exer. 2 to \(\rho_x = \varepsilon_{\pi(x)}\) . In case \(b\) ), note that \(\lambda_t\) is bounded locally almost everywhere, and reduce to showing that the mapping \(t \mapsto \pi(\lambda_t)\) is vaguely \(\mu\) -measurable (cf.~Ch. V, §3, No.~2, Prop. 4). Given a number \(\varepsilon > 0\) and a compact subset \(K\) of \(T\) , observe first that there exists in \(X\) an increasing sequence \((F_m)\) of compact sets such that the restriction of \(\pi\) to each \(F_m\) is continuous and such that, on setting \(N_m = X - F_m\) , \(\nu(N_m) \leqslant \varepsilon/4^m\) . Let \(A_m\) be the set of \(t \in K\) such that \(N_m\) is not \(\lambda_t\) -integrable or such that \(\lambda_t(N_m) > 1/2^m\) ; if \(A\) is the union of the \(A_m\) , \(\mu(A) \leqslant \varepsilon/2\) . Let \(K_1\) be a compact subset of \(K - A\) such that \(\mu(K - K_1) \leqslant \varepsilon\) and such that, for every function \(g \in \mathcal{H}(X)\) , the restriction of \(t \mapsto \int g(x) d\lambda_t(x)\) to \(K_1\) is continuous. Using Urysohn's theorem, show that for every function \(f \in \mathcal{H}(Y)\) , the restriction to \(K_1\) of \(t \mapsto \int f(\pi(x)) d\lambda_t(x)\) is the uniform limit of continuous functions.)

*5) For every \(t \in \mathbf{R}\) let \(\mathbf{f}(t)\) be the continuous function \(x \mapsto \exp(-2\pi itx)\) , which is an element of the space \(\mathcal{C} = \mathcal{C}_{\mathbf{C}}(\mathbf{R})\) of continuous complex functions on \(\mathbf{R}\) . The space \(\mathcal{C}\) is in duality with the space \(\mathcal{C}' = \mathcal{C}'_{\mathbf{C}}(\mathbf{R})\) of complex measures on \(\mathbf{R}\) with compact support; the space \(\mathcal{K} = \mathcal{K}_{\mathbf{C}}(\mathbf{R})\) of continuous complex functions with compact support may be canonically identified with a subspace of \(\mathcal{C}'\) , by identifying a function \(\varphi \in \mathcal{K}\) with the measure having density \(\varphi\) with respect to Lebesgue measure. One denotes by \(\mathcal{D}\) the subspace of \(\mathcal{K}\) formed by the indefinitely differentiable complex functions with compact support. Show that, when \(\mathcal{C}\) is equipped with the topology \(\sigma(\mathcal{C},\mathcal{C}')\) or with the topology \(\sigma(\mathcal{C},\mathcal{K})\) , \(\mathbf{f}\) is not scalarly integrable for the Lebesgue measure \(\mu\) ; however, for the topology \(\sigma(\mathcal{C},\mathcal{D})\) , \(\mathbf{f}\) is scalarly \(\mu\) -integrable and \(\int f d\mu\) is the measure \(\varepsilon_0\) ; prove that in this case, the conditions of Prop. 7 are verified.*

\begin{enumerate}
\def\labelenumi{\arabic{enumi})}
\setcounter{enumi}{5}
\item
  Let F be a distinguished (TVS, IV, §2, Exer. 4 b)) Fréchet space, F' its dual. Show that if f is a scalarly essentially integrable mapping of T into F, then \(\int f d\mu\) belongs to the bidual F'\,' of F. (Embed F in F'\,'; apply Th. 1, as well as Prop. 2 and the Lemma of the Appendix.)
\item
  \begin{enumerate}
  \def\labelenumii{\alph{enumii})}
  \tightlist
  \item
    Let \(\mathbf{F} = \overline{\mathcal{K}(\mathbf{N})}\) be the Banach space of sequences of real numbers tending to 0, \(\mathbf{F}' = \mathbf{L}^1 (\mathbf{N})\) its dual (TVS, IV, §1, Exer. 1). Let \(\mathbf{f} = (f_n)\) be the mapping of \(\mathrm{I} = [0,1]\) into \(\mathbf{F}\) such that \(f_{n}(t) = n\varphi_{\mathrm{I}_{n}}(t)\) , where \(\mathrm{I}_n = [0,1 / n]\) . Show that \(\mathbf{f}\) is scalarly integrable for the Lebesgue measure \(\mu\) on \(\mathrm{I}\) , but that \(\int \mathbf{f}d\mu\) is an element of \(\mathbf{F}''\) not in \(\mathbf{F}\) .
  \end{enumerate}
\end{enumerate}

\begin{enumerate}
\def\labelenumi{\alph{enumi})}
\setcounter{enumi}{1}
\item
  Let \(\mathbf{g}\) be the mapping of \(J = [-1, +1]\) into \(F\) defined by the conditions \(\mathbf{g}(t) = \mathbf{f}(t)\) for \(t \geqslant 0\) , \(\mathbf{g}(t) = -\mathbf{f}(-t)\) for \(t \leqslant 0\) . Then \(\mathbf{g}\) is scalarly integrable for Lebesgue measure on \(J\) , and \(\int \mathbf{g} d\mu \in F\) , but there exist functions \(h \in \mathcal{L}^{\infty}\) such that \(\int h\mathbf{g} d\mu \notin F\) .
\item
  Deduce from \(a\) ) a new proof of the fact that the topological vector space \(F\) is not isomorphic to the dual of a normed space (TVS, IV, §2, Exer. 22 c)).
\end{enumerate}

\begin{enumerate}
\def\labelenumi{\arabic{enumi})}
\setcounter{enumi}{7}
\tightlist
\item
  Let \(\mathbf{F}\) be a Hausdorff locally convex space, \(\mathbf{f}\) a scalarly locally integrable mapping of \(\mathbf{T}\) into \(\mathbf{F}\) such that, for every compact subset \(\mathbf{K}\) of \(\mathbf{T}\) , \(\int_{\mathbf{K}} \mathbf{f} d\mu \in \mathbf{F}\) .
\end{enumerate}

\begin{enumerate}
\def\labelenumi{\alph{enumi})}
\item
  Show that if \(\mathbf{f}\) is scalarly essentially integrable, and if \(\mathrm{F}\) is semi-reflexive (or if merely every Cauchy sequence for \(\sigma (\mathrm{F},\mathrm{F}^{\prime})\) converges in \(\mathrm{F}\) for this topology, when \(\mathrm{T}\) is countable at infinity), then \(\int g\mathbf{f}d\mu \in \mathrm{F}\) for every function \(g\in \mathcal{L}^{\infty}\) .
\item
  Show that if F is quasi-complete and if, for every continuous semi-norm q on F, one has \(\int^{\bullet} q(\mathbf{f}(t)) d\mu(t) < +\infty\) , then f is scalarly essentially integrable and \(\int g\mathbf{f} d\mu \in F\) for every function \(g \in L^{\infty}\) .
\end{enumerate}

(In both cases, consider the increasing directed set of compact subsets of T.)

\begin{enumerate}
\def\labelenumi{\arabic{enumi})}
\setcounter{enumi}{8}
\item
  Let G,H be two Hausdorff locally convex spaces, U a mapping of T into \(F = \mathcal{L}_{s}(G;H)\) . Assume that G is barreled, H is quasi-complete, U is \(\mu\) -measurable and \(U(K)\) is bounded for every compact subset K of T; assume, moreover, that for every continuous semi-norm q on H and every \(x \in G\) , \(\int^{\bullet} q(U(t) \cdot x) d\mu(t) < +\infty\) . Under these conditions, show that U is scalarly essentially \(\mu\) -integrable and that \(\int U(t) d\mu(t) \in F\) (make use of Exercise 8 b)).
\item
  Let \(G, H\) be two Fréchet spaces, \(G_0\) (resp. \(H_0\) ) a dense subset of \(G\) (resp. \(H\) ), and \(t \mapsto \Phi_t\) a mapping of \(T\) into the space \(F = \mathcal{B}(G, H)\) of continuous bilinear forms on \(G \times H\) , equipped with the topology of pointwise convergence. Assume that: \(1^\circ\) for every pair \((a, b) \in G_0 \times H_0\) , \(t \mapsto \Phi_t(a, b)\) is essentially \(\mu\) -integrable; \(2^\circ\) if, for every pair of bounded subsets \(A \subset G\) , \(B \subset H\) , one sets \(q_{A,B}(\Phi) = \sup_{(x, y) \in A \times B} |\Phi(x, y)|\) for
\end{enumerate}

every \(\Phi \in \mathcal{B}(\mathrm{G},\mathrm{H})\) , then \(\int^{*}q_{\mathrm{A},\mathrm{B}}(\Phi_t)d\mu (t) < + \infty\) . Under these conditions, show that \(t\mapsto \Phi_t\) is scalarly essentially \(\mu\) -integrable and \(\int \Phi_t d\mu (t)\in \mathrm{F}\) (using Lebesgue's theorem, reduce to applying Cor. 1 of Th. 1). Special case where \(\mathrm{H} = \mathbf{R}\) .

\begin{enumerate}
\def\labelenumi{\arabic{enumi})}
\setcounter{enumi}{10}
\item
  Let \(\mu\) be Lebesgue measure on \(T = [0,1]\) , \(F\) a Hilbert space having a countable orthonormal basis \((\mathbf{e}_n)\) . Let \(\mathbf{f}\) be the mapping of \(T\) into \(F\) such that \(\mathbf{f}(0) = 0\) and \(\mathbf{f}(t) = 2^n \mathbf{e}_n / n\) on the interval \([2^{-n-1}, 2^{-n}]\) ( \(n \geqslant 0\) ). Show that \(\mathbf{f}\) is \(\mu\) -measurable, scalarly \(\mu\) -integrable, and that \(\int \mathbf{f} d\mu \in F\) , but that \(\int^* |\mathbf{f}| d\mu = +\infty\) .
\item
  Let \(\mu\) be Lebesgue measure on \(T = [0,1]\) , and let \(F\) be a Hilbert space having an orthonormal basis \((\mathbf{e}_t)_{t \in T}\) equipotent to \(T\) .
\end{enumerate}

\begin{enumerate}
\def\labelenumi{\alph{enumi})}
\item
  Let \(\mathbf{f}\) be the mapping \(t \mapsto \mathbf{e}_t\) of \(\mathrm{T}\) into \(\mathrm{F}\) . Show that \(\mathbf{f}\) is scalarly \(\mu\) -negligible, but is not \(\mu\) -measurable when \(\mathrm{F}\) is equipped with the weak topology \(\sigma(\mathrm{F}, \mathrm{F}')\) (nor, a fortiori, when \(\mathrm{F}\) is equipped with its original topology).
\item
  Let A be a non-measurable set in T (Ch. IV, §4, Exer. 8); show that the function \(\mathbf{g} = \mathbf{f}\varphi_{\mathrm{A}}\) is scalarly \(\mu\) -negligible, but that the numerical function \(|\mathbf{g}|\) is not \(\mu\) -measurable.
\end{enumerate}

\begin{enumerate}
\def\labelenumi{\arabic{enumi})}
\setcounter{enumi}{12}
\item
  Let F be a Hausdorff locally convex space, \(F'\) its dual, q a lower semi-continuous semi-norm on F. Show that if f is a mapping of T into F, \(\mu\) -measurable for the topology \(\sigma(\mathrm{F},\mathrm{F}')\) , then the numerical function \(q \circ f\) is \(\mu\) -measurable.
\item
  Let \(\mu\) be Lebesgue measure on T = {[}0,1{]}; denote by F the vector space over R of \(\mu\) -measurable finite numerical functions on T, equipped with the topology of pointwise convergence, which makes it a Hausdorff locally convex space.
\end{enumerate}

\begin{enumerate}
\def\labelenumi{\alph{enumi})}
\item
  Show that there exists in F a countable dense subset (consider a dense sequence in the Banach space \(\mathcal{C}(T)\) of continuous numerical functions on T).
\item
  For every \(t \in T\) , let \(\mathbf{f}(t)\) be the element of the dual \(F'\) of F defined by \(\langle\mathbf{f}(t), z\rangle = z(t)\) for all \(t \in T\) . Show that, when \(F'\) is equipped with the topology \(\sigma(\mathrm{F}', \mathrm{F})\) , f is scalarly \(\mu\) -measurable but is not \(\mu\) -measurable (cf.~Prop. 13).
\end{enumerate}

\begin{enumerate}
\def\labelenumi{\arabic{enumi})}
\setcounter{enumi}{14}
\item
  Let F be a metrizable locally convex space, f a scalarly \(\mu\) -measurable mapping of T into F such that, for every compact subset K of T, there exists a countable subset H of F such that \(\mathbf{f}(t) \in \overline{\mathbf{H}}\) for almost every \(t \in K\) ; show that under these conditions, f is \(\mu\) -measurable for the original topology of F. (Embed F in a countable product of normed spaces, after having reduced to the case that F is separable.)
\item
  Let \(\mathbf{F}\) be a Banach space, \(\mathbf{F}'\) its dual. For every \(p\) such that \(1 \leqslant p \leqslant +\infty\) , denote by \(\Lambda_{\mathbf{F}'}^p(\mathrm{T},\mu)\) (or simply by \(\Lambda_{\mathbf{F}'}^p\) ) the set of mappings \(\mathbf{f}\) of \(\mathrm{T}\) into \(\mathbf{F}'\) such that, for every \(\mathbf{z} \in \mathbf{F}\) , the numerical function \(\langle \mathbf{z},\mathbf{f}\rangle\) belongs to \(\mathcal{L}^p(\mathrm{T},\mu)\) . One has \(\Lambda_{\mathbf{F}'}^p \supset \mathcal{L}_{\mathbf{F}'}^p\) , and \(\Lambda_{\mathbf{F}'}^1\) is the space of scalarly essentially \(\mu\) -integrable functions with values in \(\mathbf{F}'\) ( \(\mathbf{F}'\) being equipped with \(\sigma(\mathbf{F}',\mathbf{F})\) ). One knows that, denoting by \(\theta(\mathbf{z})\) the class of \(\langle \mathbf{z},\mathbf{f}\rangle\) in \(\mathrm{L}^p(\mu)\) , the mapping \(\mathbf{z} \mapsto \theta(\mathbf{z})\) is continuous (No.~4, Lemma 2); let \(M_p(\mathbf{f})\) be its norm; this is a semi-norm on the space \(\Lambda_{\mathbf{F}'}^p\) . In order that \(M_p(\mathbf{f}) = 0\) , it is necessary and sufficient that \(\mathbf{f}\) be scalarly locally negligible.
\end{enumerate}

\begin{enumerate}
\def\labelenumi{\alph{enumi})}
\item
  In order that \(\mathbf{f} \in \Lambda_{\mathbf{F}'}^p\) , it is necessary and sufficient that for every numerical function \(g \in \mathcal{L}^q\) , the function \(g\mathbf{f}\) be scalarly essentially \(\mu\) -integrable; one then has \(\mathrm{M}_1(g\mathbf{f}) \leqslant \mathrm{M}_p(\mathbf{f})\mathrm{N}_q(g)\) .
\item
  In the space \(\Lambda_{\mathbf{F}'}^1\) , show that the semi-norm \(\mathbf{M}_1\) is equivalent to the semi-norm \(\mathrm{M}_1'(\mathbf{f}) = \sup |\int_{\mathbf{A}} \mathbf{f} d\mu|\) , where A runs over the set of measurable subsets of T; more precisely, \(\mathrm{M}_1' \leqslant \mathrm{M}_1 \leqslant 2\mathrm{M}_1'\) . Deduce from this that \(\mathrm{M}_p\) is a semi-norm equivalent to \(\mathrm{M}_p'(\mathbf{f}) = \sup_{\mathrm{N}_q(g) \leqslant 1} \mathrm{M}_1'(g\mathbf{f})\) .
\item
  Take for F a Hilbert space having a countable orthonormal basis, and for \(\mu\) the Lebesgue measure on \(T = [0,1]\) . Show that there exists a mapping of T into F, not integrable but scalarly integrable, that is the limit for the semi-norm \(M_1\) of a sequence of measurable step functions (cf.~Exer. 11); from this, deduce that on the space \(\mathcal{L}_{\mathrm{F}}^1\) , of \(\mu\) -integrable functions, the topology defined by the semi-norm \(M_1\) is strictly coarser than that defined by \(N_1\) . Analogous result for the semi-norms \(M_p\) and \(N sub p\) , for \(1 \leqslant p < +\infty\) .
\end{enumerate}

\(d\) ) Show that on the space \(\mathcal{L}_{\mathrm{F}^{\prime}}^{\infty}\) , \(\mathrm{M}_{\infty} = \mathrm{N}_{\infty}\)

\begin{enumerate}
\def\labelenumi{\alph{enumi})}
\setcounter{enumi}{4}
\tightlist
\item
  Take for F a Hilbert space having a countable orthonormal basis, arranged in a double sequence \((\mathbf{e}_{mn})\) . Let \(\mu\) be Lebesgue measure on \(T = [0,1]\) . Let \(\mathbf{u}_m\) be the mapping of T into F such that \(\mathbf{u}_m(t) = \mathbf{e}_{mn}\) for \((n - 1)2^{-m} \leqslant t < n2^{-m}\) and for \(1 \leqslant n \leqslant 2^m\) , and \(\mathbf{u}_m(1) = 0\) ; set \(\mathbf{f}_n = \sum_{i=0}^{n} \mathbf{u}_i\) . Show that, for \(1 \leqslant p < +\infty\) , \((\mathbf{f}_n)\) is a Cauchy sequence in \(\Lambda_{\mathrm{F}'}^p\) but does not converge to any function in \(\Lambda_{\mathrm{F}'}^p\) .
\end{enumerate}

\begin{enumerate}
\def\labelenumi{\alph{enumi})}
\setcounter{enumi}{5}
\item
  Let \((\mathbf{f}_{n})\) be a Cauchy sequence in the space \(\Lambda_{F'}^{p}\) ( \(1 \leqslant p \leqslant +\infty\) ); assume that there exists a mapping f of T into \(F'\) such that, for every \(z \in F\) , the sequence of functions \(\langle z, f_{n} \rangle\) converges in measure to \(\langle z, f \rangle\) (Ch. IV, §5, No.~11). Show that \(f \in \Lambda_{F'}^{p}\) and that the sequence \((\mathbf{f}_{n})\) has limit f in \(\Lambda_{F'}^{p}\) .
\item
  Take for F the Banach space \(\mathrm{L}^{1}(\mathbf{N})\) of absolutely convergent series; its dual is \(\mathbf{F}^{\prime} = \mathrm{L}^{\infty}(\mathbf{N})\) (TVS, IV, §1, Exer. 1). Let \(\mu\) be Lebesgue measure on T = {[}0,1{]}. Let f be the mapping of T into \(F^{\prime}\) such that \(\mathbf{f}(0) = 0\) and, for \(2^{-n-1} < t \leqslant 2^{-n}\) , \(\mathbf{f}(t)\) is the sequence all of whose terms are zero, except for the term of index n, equal to \(2^{n/p}\) for \(1 \leqslant p < +\infty\) . Show that f is measurable for the strong topology on \(F^{\prime}\) , and belongs to \(\Lambda_{\mathrm{F}'}^p\) , but that there does not exist any sequence \((\mathbf{f}_n)\) of step functions that tends to \(\mathbf{f}\) in the space \(\Lambda_{\mathrm{F}'}^p\) .
\end{enumerate}

¶ 17) Let \(\mu\) be Lebesgue measure on \(T = [0,1]\) , \(F\) the Banach space \(L^1(\mathbf{N})\) , \(F' = L^\infty(\mathbf{N})\) its dual. For every \(t \in [0,1[\) , let \(t = \sum_{n=0}^\infty \xi_n 2^{-n}\) be the dyadic expansion of \(t\) ; denote by \(f(t)\) the sequence \((\xi_n) \in F'\) , and set \(f(1) = 0\) .

\begin{enumerate}
\def\labelenumi{\alph{enumi})}
\item
  Show that, for \(1 \leqslant p \leqslant +\infty\) , the function \(\mathbf{f}\) belongs to the space \(\Lambda_{\mathrm{F}'}^p\) (cf.~Exer. 16).
\item
  Show that \(\mathbf{f}\) is measurable for the topology \(\sigma(\mathrm{F}',\mathrm{F})\) and that, for this topology, there exists a sequence of step functions that converges almost everywhere to \(\mathbf{f}\) .
\item
  Show that \(\mathbf{f}\) is not measurable for the strong topology on \(\mathbf{F}'\) (observe that \(|\mathbf{f}(t) - \mathbf{f}(t')| = 1\) if \(0 \leqslant t < t' < 1\) ).
\item
  Show that, in the space \(\Lambda_{\mathrm{F}'}^p\) ( \(1 \leqslant p < +\infty\) ), \(\mathbf{f}\) is not the limit of a sequence of step functions. (Reduce to considering only step functions that are linear combinations (with coefficients in \(\mathbf{F}'\) ) of characteristic functions of intervals whose endpoints are of the form \(k/2^n\) ; if \(\mathbf{g}\) is such a step function, show that, in the notations of Exer. 16 b), \(\mathrm{M}_1'(\mathbf{f} - \mathbf{g}) \geqslant 1/4\) .)
\item
  Show that there does not exist any mapping \(\mathbf{h}\) of \(\mathrm{T}\) into \(\mathrm{F}'\) , measurable for the strong topology, such that \(\mathbf{f} - \mathbf{h}\) is scalarly negligible. (Observe that, for such a mapping, one would necessarily have \(\mathrm{N}_{\infty}(\mathbf{h}) \leqslant 1\) ; then obtain a contradiction with the result of \(d\) .)
\end{enumerate}

\begin{enumerate}
\def\labelenumi{\arabic{enumi})}
\setcounter{enumi}{17}
\tightlist
\item
  Show that Prop. 8 continues to hold when the condition that \(\mathbf{f}\) be \(\mu\) -measurable (for the original topology of \(\mathbf{F}\) ) is replaced by the condition that \(\mathbf{f}\) be \(\mu\) -measurable for the weakened topology \(\sigma(\mathrm{F},\mathrm{F}')\) . (Make use of Exer. 19 c) of Ch. IV, §4.)
\end{enumerate}

¶ 19) Let f be a scalarly essentially μ-integrable mapping of T into F. For every \(z' \in F'\) , denote by \(u_{\mathbf{f}}(\mathbf{z}')\) the class in \(L^{1}(\mu)\) of the function \(\langle f, z' \rangle\) . We shall say that f is scalarly well integrable if, for every function \(g \in L^{\infty}(\mu)\) , \(\int g f d\mu \in F\) .

\begin{enumerate}
\def\labelenumi{\alph{enumi})}
\item
  For \(\mathbf{f}\) to be scalarly well integrable, it is necessary and sufficient that \(u_{\mathbf{f}}\) be continuous for the topologies \(\sigma(\mathrm{F}',\mathrm{F})\) and \(\sigma(\mathrm{L}^1,\mathrm{L}^\infty)\) . When this is so, the image under \(u_{\mathbf{f}}\) of every equicontinuous subset of \(\mathrm{F}'\) is a relatively weakly compact subset of \(\mathrm{L}^1\) . In particular: \(\alpha\) for every equicontinuous subset \(\mathrm{H}'\) of \(\mathrm{F}'\) and every \(\varepsilon >0\) , there exists a compact subset \(\mathrm{L}\) of \(\mathrm{T}\) such that \(\int_{\mathrm{T}-\mathrm{L}}|\langle\mathbf{f},\mathbf{z}'\rangle|d\mu\leqslant\varepsilon\) for all \(\mathbf{z}'\in\mathrm{H}'\) ; \(\beta\) for every equicontinuous subset \(\mathrm{H}'\) of \(\mathrm{F}'\) and every \(\varepsilon>0\) , there exists an \(\eta>0\) such that, for every open set \(\mathrm{U}\subset\mathrm{T}\) of measure \(\mu(\mathrm{U})\leqslant\eta\) , one has \(\int_{\mathrm{U}}|\langle\mathbf{f},\mathbf{z}'\rangle|d\mu\leqslant\varepsilon\) for all \(\mathbf{z}'\in\mathrm{H}'\) (Ch. V, §5, Exer. 15).
\item
  Assume that the conditions \(\alpha\) and \(\beta\) of \(a\) ) are verified and, moreover, that there exists a subspace \(\mathrm{H} \subset \mathrm{L}^{\infty}\) , dense for the topology \(\sigma(\mathrm{L}^{\infty}, \mathrm{L}^{1})\) , such that for every function \(g \in \mathcal{L}^{\infty}\) whose class is in \(\mathrm{H}\) , one has \(\int g\mathbf{f}d\mu \in \mathrm{F}\) . Show that under these conditions, if in addition \(\mathrm{F}\) is assumed to be quasi-complete, then \(\mathbf{f}\) is scalarly well integrable. (First reduce to the case that \(\mathrm{F}\) is complete (for its original topology), and observe that \(\dot{g} \mapsto \int g\mathbf{f}d\mu\) is a continuous linear mapping of \(\mathrm{L}^{\infty}\) into \(\mathrm{F}'*\) , when \(\mathrm{L}^{\infty}\) is equipped with the topology \(\tau(\mathrm{L}^{\infty}, \mathrm{L}^{1})\) and \(\mathrm{F}'*\) with the topology of uniform convergence on the equicontinuous subsets of \(\mathrm{F}'\) .)
\item
  Suppose that \(\mathbf{f}\) is scalarly well integrable. Show that for every Cauchy sequence \((\mathbf{z}_n')\) for the topology \(\sigma(\mathrm{F}',\mathrm{F})\) , the sequence \((u_{\mathbf{f}}(\mathbf{z}_n'))\) is strongly convergent in \(\mathrm{L}^1\) (use Exer. 25 c) of Ch. IV, §5 and Exer. 16 b) of Ch. V, §5). From this, deduce that if A is the image of the unit ball of \(\mathrm{L}^{\infty}\) under the mapping \(\dot{g} \mapsto \int g\mathbf{f}d\mu\) , then every Cauchy sequence \((\mathbf{z}_n')\) for the topology \(\sigma(\mathrm{F}',\mathrm{F})\) is a Cauchy sequence for the topology of uniform convergence in A.
\end{enumerate}

\(d)\) Deduce from \(c)\) that if \(\mathbf{f}\) is scalarly well integrable, then the image under \(u_{\mathbf{f}}\) of every subset of \(\mathrm{F}'\) that is equicontinuous and metrizable for \(\sigma (\mathrm{F}',\mathrm{F})\) , is relatively strongly compact in \(L^{1}\) . In particular, if F is the dual \(G'\) of a Hausdorff locally convex space G, if F is equipped with a topology compatible with the duality between F and G, and if there exists in \(G'\) a countable strongly dense subset, then the image under \(u_{f}\) of every bounded subset of \(F' = G\) is relatively strongly compact in \(L^{1}\) (embed G in its bidual).

\begin{enumerate}
\def\labelenumi{\arabic{enumi})}
\setcounter{enumi}{19}
\tightlist
\item
  Let F be a Hausdorff locally convex space, \(F'\) its dual, \(S'\) the set of subsets \(A' \subset F'\) such that every sequence of points of \(A'\) has a subsequence convergent for \(\sigma(\mathrm{F}',\mathrm{F})\) . Show that, for a bounded subset A of F, the following conditions are equivalent: \(\alpha\) A is precompact for the \(S'\) -topology; \(\beta\) every sequence convergent for \(\sigma(\mathrm{F}',\mathrm{F})\) converges uniformly on A. (Let S be the set formed by the finite subsets of F and of A; observe that \(\beta\) ) implies that every set \(A' \in S'\) is precompact for the S-topology, using Exer. 6 of GT, II, §4. Then apply Exer. 4 of TVS, IV, §1.)
\end{enumerate}

¶ 21) Let f be a scalarly well integrable mapping of T into a quasi-complete Hausdorff locally convex space F.

\begin{enumerate}
\def\labelenumi{\alph{enumi})}
\item
  Show that the image under \(\dot{g} \mapsto \int g\mathbf{f}  d\mu\) of the unit ball of \(\mathrm{L}^{\infty}\) is a relatively compact subset of \(\mathbf{F}\) for the original topology, in each of the following cases: \(1^{\circ}\) there exists a countable dense subset of \(\mathbf{F}\) ; \(2^{\circ}\) \(\mathbf{F}\) is the dual \(\mathbf{G}'\) of a space \(\mathbf{G}\) that is either metrizable or is the strict direct limit of a sequence of metrizable spaces, and \(\mathbf{F}\) is equipped with the topology \(\tau (\mathbf{G}',\mathbf{G})\) . (Make use of Exer. 19 c) and Exer. 20; in case \(2^{\circ}\) , use Šmulian's theorem (TVS, IV, §5, Exer. 2 c)).
\item
  Show that the conclusion of a) remains valid when, with no new hypothesis on F, one assumes f to be measurable. (Embedding F in a product of Banach spaces, reduce to the case that F is a Banach space; then, using Exer. 19 a), reduce to the case that T is compact; then apply case \(1^{\circ}\) of a).
\end{enumerate}

\begin{enumerate}
\def\labelenumi{\arabic{enumi})}
\setcounter{enumi}{21}
\tightlist
\item
  In a Hausdorff locally convex space \(\mathbf{F}\) , let \((\mathbf{z}_{\alpha})_{\alpha \in \mathbf{A}}\) be a family such that the mapping \(\alpha \mapsto \mathbf{z}_{\alpha}\) is scalarly well integrable (Exer. 19) when \(\mathbf{A}\) is equipped with the measure defined by mass \(+1\) at each point.
\end{enumerate}

\begin{enumerate}
\def\labelenumi{\alph{enumi})}
\item
  Show that the family \((\mathbf{z}_{\alpha})\) is summable for every topology T compatible with the duality between F and \(F'\) (make use of Exer. 19 a)). Converse when F is quasi-complete for T.
\item
  Suppose that \(F = G'\) , \(F' = G\) , where G is a Hausdorff locally convex space such that there exists in \(G'\) a countable strongly dense subset. Show that \((\mathbf{z}_{\alpha})\) is then summable for the strong topology on \(G'\) (which is not necessarily compatible with the duality between \(G'\) and G (cf.~Exer. 19d)). Show that this result does not extend to the case that \(G = L^{1}(\mathbf{N})\) , \(G' = L^{\infty}(\mathbf{N})\) .
\end{enumerate}

\begin{enumerate}
\def\labelenumi{\arabic{enumi})}
\setcounter{enumi}{22}
\tightlist
\item
  Let \(\mu\) be a positive measure on \(T\) , carried by a countable union of compact sets.
\end{enumerate}

\begin{enumerate}
\def\labelenumi{\alph{enumi})}
\tightlist
\item
  Let H be a set of \(\mu\) -measurable numerical functions, such that for every \(t \in \mathbf{T}\) , \(\sup_{f \in \mathrm{H}} f(t) < +\infty\) . Show that there exists a finite \(\mu\) -measurable function \(g\) such that
\end{enumerate}

\(f(t) \leqslant g(t)\) almost everywhere for every function \(f \in \mathrm{H}\) . (Reduce to the case that \(\mu\) is bounded and the functions in \(\mathbf{H}\) take their values in \([-1, +1]\) , by means of an increasing homeomorphism of \(\overline{\mathbf{R}}\) onto this interval. Then consider the supremum \(\widetilde{g}\) in \(\mathrm{L}^1(\mu)\) of the set of classes of the functions in \(\mathbf{H}\) (Ch. IV, §3, No.~6, Prop. 14) and observe that there exists a sequence of functions in \(\mathbf{H}\) that converges almost everywhere to \(g\) .)

\begin{enumerate}
\def\labelenumi{\alph{enumi})}
\setcounter{enumi}{1}
\item
  Let \(\mathbf{f}\) be a scalarly \(\mu\) -measurable mapping of \(\mathrm{T}\) into a Hausdorff locally convex space \(\mathrm{F}\) . For every subset \(\mathrm{B}'\) of \(\mathrm{F}'\) , bounded for \(\sigma(\mathrm{F}',\mathrm{F})\) , show that there exists a measurable and finite numerical function \(g_{\mathrm{B}'}\) such that, for every \(\mathbf{z}' \in \mathrm{B}'\) , \(|\langle \mathbf{f}(t),\mathbf{z}' \rangle| \leqslant g_{\mathrm{B}'}(t)\) almost everywhere.
\item
  Let \(\mathbf{F}\) be a Fréchet space, \(\mathbf{f}\) a scalarly \(\mu\) -measurable mapping of \(\mathbf{T}\) into \(\mathbf{F}\) . Show that for every compact subset \(\mathbf{K}\) of \(\mathbf{T}\) and every \(\varepsilon > 0\) , there exists a compact set \(\mathrm{K}' \subset \mathrm{K}\) such that \(\mu(\mathrm{K} - \mathrm{K}') \leqslant \varepsilon\) and the restriction of \(\mathbf{f}\) to \(\mathrm{K}'\) is bounded (make use of \(b\) )).
\end{enumerate}

\begin{enumerate}
\def\labelenumi{\arabic{enumi})}
\setcounter{enumi}{23}
\item
  Let F be a complete Hausdorff locally convex space, containing a countable dense subset. Let f be a mapping of T into F, scalarly \(\mu\) -measurable and scalarly essentially bounded; show that for every function \(g \in L^{1}\) , gf is scalarly integrable and \(\int g f d\mu \in F\) . (Make use of Cor. 1 of Th. 2 of TVS, III, §3, No.~6, on observing that the equicontinuous subsets of \(F'\) are metrizable for \(\sigma(F', F)\) , and applying Lebesgue's theorem.)
\item
  Let F be a separable Fréchet space in which every Cauchy sequence for \(\sigma(F, F')\) is convergent for this topology (for example a space \(L^1(T', \nu)\) , where \(T'\) has a countable base (Ch. V, §5, Exer. 16 b))). Show that every scalarly essentially integrable mapping f of T into F is scalarly well integrable. (First prove that for every function \(g \in \mathcal{L}^\infty(\mu)\) with compact support, one has \(\int g\mathbf{f} d\mu \in F\) , on observing that f is \(\mu\) -measurable (No.~5, Prop. 12), and on applying Prop. 8 of No.~2 and the hypothesis. Deduce from this, with the help of the hypothesis, that \(\int_A g\mathbf{f} d\mu \in F\) for every function \(g \in \mathcal{L}^\infty\) and every set A that is a countable union of compact sets. Finally, with the notations of Exer. 19, prove that the image under \(u_{\mathbf{f}}\) of every equicontinuous subset of \(F'\) is relatively weakly compact in \(L^1\) , using Eberlein's theorem (TVS, IV, §5, No.~3, Th. 1) and the fact that the equicontinuous subsets of \(F'\) are metrizable for \(\sigma(F', F)\) .)
\item
  Let \((f_n)\) be a sequence of \(\mu\) -integrable functions defined on \(\mathbf{T}\) , such that for almost every \(t \in \mathbf{T}\) , \(\sup_{n} |f_n(t)| < +\infty\) and \(\sup_{n} \int |f_n(t)| d\mu(t) = +\infty\) . Show that there exists a sequence \((c_n)\) of scalars such that \(\sum_{n} |c_n| < +\infty\) and such that \(\sum_{n} c_n f_n\) is not \(\mu\) -integrable (apply the Gelfand-Dunford theorem to the space \(\mathbf{L}^1(\mathbf{N}) = \mathbf{F}\) .)
\item
  Let E be a separable real Banach space, \(G = E'\) its dual, \(G'\) a subspace of E that is dense in E, distinct from E, and barreled (TVS, III, §4, Exer. 4). Show that if G is equipped with the topology \(\sigma(G, G')\) , and its dual \(G' = \mathcal{L}(G; \mathbf{R})\) with the topology of compact convergence, then \(G'\) is not quasi-complete, and there exist a sequence \((x_n)\) in \(G'\) tending to 0 and a sequence \((\lambda_n)\) of numbers \textgreater0 and with finite sum, such that the series with general term \(\lambda_n x_n\) is not convergent in \(G'\) . (Observe first that every subset of G bounded for \(\sigma(G, G')\) is bounded for the norm on G. Consider a point \(a \in E\) not belonging to \(G'\) , and a sequence \((a_n)\) of points of \(G'\) that converges to a in E and is such that \(\|a_{n+1} - a_n\| \leqslant 1/2^n\) .)
\end{enumerate}

\setcounter{exercisegroup}{2}
\edef\CurrentExerciseGroup{\arabic{exercisegroup}}
\section*{\texorpdfstring{\S\CurrentExerciseGroup}{Section \CurrentExerciseGroup}}
\phantomsection
\addcontentsline{toc}{section}{\texorpdfstring{Exercises for \S\CurrentExerciseGroup}{Exercises for Section \CurrentExerciseGroup}}

\begin{enumerate}
\def\labelenumi{\arabic{enumi})}
\item
  Show that, for a numerical function \(f\) on \(T\) to be essentially integrable for every positive measure on \(T\) , it is necessary and sufficient that it be bounded, have compact support, and be measurable for every positive measure on \(T\) .
\item
  Let \(G\) be a barreled space, \(H\) a semi-reflexive space, and \(F\) the space \(\mathcal{L}_s(G;H)\) . Let \(m\) be a vectorial measure on \(T\) , with values in \(F\) . Show that for every numerical function \(f\) essentially integrable for \(m\) , \(\int fd\mathbf{m} \in F\) . (Observe that \(F\) is semi-reflexive.) The case where \(H = R\) .
\item
  Let \(\mathbf{m}\) be a vectorial measure on \(\mathrm{T}\) , with values in \(\mathrm{F}\) . Show that for every numerical function \(f\) essentially integrable for \(\mathbf{m}\) , one has \(\int f dm \in \mathrm{F}''\) in each of the following two cases: \(a)\) \(\mathrm{F}\) is distinguished (TVS, IV, §2, Exer. 4); \(b)\) \(\mathrm{F}\) is a Fréchet space and \(\mathrm{T}\) is countable at infinity. (In both cases, embed \(\mathrm{F}\) in \(\mathrm{F}''\) and apply the method of Cor. 1 of Prop. 3; in case \(b\) ), use the Cor. of Prop. 3 of TVS, IV, §3, No.~3.)
\item
  Let \(\mathbf{m}\) be a vectorial measure on \(\mathrm{T}\) , with values in \(\mathrm{F}\) . Show that for every numerical function \(f \geqslant 0\) that is lower semi-continuous and essentially integrable for \(\mathbf{m}\) , one has \(\int f d\mathbf{m} \in \mathrm{F}''\) (use Th. 1 of Ch. IV, §1, No.~1).
\item
  Let F be a semi-reflexive locally convex space, m a measure on T with values in F. Show that for every numerical function g locally integrable for m, the mapping \(f \mapsto \int fg dm\) of \(\mathcal{K}(T)\) into F is a measure on T, denoted \(g \cdot m\) ; in order that a numerical function h be essentially integrable for \(g \cdot m\) , it is necessary and sufficient that gh be essentially integrable for m, in which case \(\int h d(g \cdot m) = \int hg dm\) .
\item
  Let F be a semi-reflexive locally convex space, equipped with an algebra structure over R, such that the mapping \((u,v)\mapsto uv\) of \(F\times F\) into F is separately continuous. Let m be a measure on T, with values in F, such that \(\mathbf{m}(fg)=\mathbf{m}(f)\mathbf{m}(g)\) for f and g in \(\mathcal{K}(\mathrm{T})\) . Show that if the numerical functions f,g are such that f,g, and fg belong to \(\mathcal{L}(\mathbf{m})\) , then \(\int fg dm=(\int fd\mathbf{m})(\int gd\mathbf{m})\) . (First treat the case that \(g\in\mathcal{K}(\mathrm{T})\) , by considering the two measures \(f\mapsto\mathbf{m}(fg)\) and \(f\mapsto\mathbf{m}(f)\mathbf{m}(g)\) ; pass to the general case by proceeding similarly and making use of Exer. 5.) The case where \(F=\mathcal{L}(\mathrm{G})\) , G being reflexive and F equipped with the topology of pointwise convergence or the topology of weak pointwise convergence.
\item
  Let \(\mu\) be a positive measure on T, m a measure on T with values in F, having base \(\mu\) and density f with respect to \(\mu\) . Let \(q\) be a lower semi-continuous semi-norm on F. Show that if, for every compact subset K of T, \(\int_{K}^{*}(q \circ f) d\mu\) is finite, then, for every \(\mu\) -measurable function \(h \geqslant 0\) , the inequality \(\int^{\bullet} h dq(m) \leqslant \int^{\bullet} (q \circ f) h d\mu\) holds. (Use Lemma 1 of Ch. V, 1st edn., §2, No.~2.)²
\item
  Let \(\mathbf{m}\) be a vectorial measure on \(\mathbf{T}\) with values in \(\mathbf{F}\) , \(q\) -majorizable for a lower semi-continuous semi-norm \(q\) on \(\mathbf{F}\) . Show that for every function \(f \geqslant 0\) of \(\mathcal{H}(\mathbf{T})\) ,
\end{enumerate}

\[
\int f d (q \circ \mathbf {m}) = \sup \sum_ {i} q (\mathbf {m} (f _ {i})),
\]

where \((f_i)\) runs over the set of all finite sequences of elements of \(\mathcal{H}(\mathrm{T})\) such that \(\sum_{i}|f_i| \leqslant f\) . (Argue as in Th. 1 of Ch. II, §2, No.~2.)

\begin{enumerate}
\def\labelenumi{\arabic{enumi})}
\setcounter{enumi}{8}
\tightlist
\item
  Let T be a locally compact space countable at infinity, F a Hausdorff barreled space containing a countable dense subset, m a vectorial measure on T with values in the weak dual F' of F. Show that there exists a positive measure \(\mu\) on T such that m is scalarly of base \(\mu\) . (Make use of Prop. 12 of Ch. V, §5, No.~6, and condition 5) of Cor. 5 of the Lebesgue--Nikodym theorem, Ch. V, §5, No.~5, Th. 2.)
\end{enumerate}

¶ 10) Let K be a compact space. For every Hausdorff quotient space \(K_{1}\) of K, one identifies the Banach space \(\mathcal{C}(K_{1})\) with a closed subspace of \(\mathcal{C}(K)\) by means of the injection \(f \mapsto f \circ \pi\) , where \(\pi\) is the canonical mapping of K onto \(K_{1}\) .

\begin{enumerate}
\def\labelenumi{\alph{enumi})}
\item
  Let \(u\) be a continuous linear mapping of \(\mathcal{C}(\mathrm{K})\) into a quasi-complete Hausdorff locally convex space \(\mathrm{F}\) . In order that \(u\) transform the unit ball of \(\mathcal{C}(\mathrm{K})\) into a subset of \(\mathrm{F}\) that is relatively compact for \(\sigma(\mathrm{F},\mathrm{F}')\) , it suffices that, for every metrizable quotient space \(\mathrm{K}_1\) of \(\mathrm{K}\) , the restriction of \(u\) to \(\mathcal{C}(\mathrm{K}_1)\) have that property. (Make use of Eberlein's theorem (TVS, IV, §5, No.~3, Th. 1), and observe that every sequence \((f_n)\) in \(\mathcal{C}(\mathrm{K})\) is contained in a subspace \(\mathcal{C}(\mathrm{K}_1)\) , where \(\mathrm{K}_1\) is a metrizable quotient space of \(\mathrm{K}\) ; for that, one considers the continuous mapping \(x \mapsto (f_n(x))\) of \(\mathrm{K}\) into \(\mathbf{R}^{\mathbf{N}}\) .)
\item
  For every Hausdorff quotient space \(K_{1}\) of K, the transpose of the canonical injection of \(\mathcal{C}(K_{1})\) into \(\mathcal{C}(K)\) is the mapping \(\mu\mapsto\pi(\mu)\) of \(\mathcal{M}(K)\) into \(\mathcal{M}(K_{1})\) . Denote by \((\mathcal{M}(\mathrm{K}))'\) the dual of the Banach space \(\mathcal{M}(\mathrm{K})\) (the bidual of \(\mathcal{C}(\mathrm{K})\) ). For a subset A of \(\mathcal{M}(\mathrm{K})\) to be relatively compact for \(\sigma(\mathcal{M}(\mathrm{K}),(\mathcal{M}(\mathrm{K}))')\) , it is necessary and sufficient that, for every metrizable quotient space \(K_{1}\) of K, the canonical image of A in \(\mathcal{M}(\mathrm{K}_{1})\) be relatively compact for \(\sigma(\mathcal{M}(\mathrm{K}_{1}),(\mathcal{M}(\mathrm{K}_{1}))')\) . (Observe that A is strongly bounded; one can then suppose that A is convex, balanced and closed for the vague topology \(\sigma(\mathcal{M}(\mathrm{K}),\mathcal{C}(\mathrm{K}))\) (apply in \(\mathcal{M}(\mathrm{K}_{1})\) Exer. 10 of Ch. IV, §7); A is therefore the image of the unit ball of the dual \(F'\) of a Banach space F, under the transpose of a continuous linear mapping u of \(\mathcal{C}(\mathrm{K})\) into F (consider the polar of A in \(\mathcal{C}(\mathrm{K})\) ); then apply a).)
\end{enumerate}

¶ 11) Let F, G be two Hausdorff locally convex spaces, G being assumed quasi-complete; let u be a continuous linear mapping of F into G; its transpose \(^{t}u\) being strongly continuous, the transpose \(^{t}(^{t}u)\) is a strongly and weakly continuous linear mapping of F'\,' into G'\,'. Show that the following conditions are equivalent:

\(\alpha)\) \(u\) transforms every bounded subset of \(\mathbf{F}\) into a relatively compact subset of \(\mathbf{G}\) for \(\sigma(\mathbf{G},\mathbf{G}')\) .

\(\beta)^{t}(^{t}u)\) maps \(\mathbf{F}^{\prime \prime}\) into G.

\(\gamma)\ ^{t}u\) is continuous for the topologies \(\sigma(\mathrm{G}^{\prime},\mathrm{G})\) and \(\sigma(\mathrm{F}^{\prime},\mathrm{F}^{\prime\prime})\) .

\(\delta)^{t}u\) transforms every equicontinuous subset of \(\mathrm{G}'\) into a set relatively compact for \(\sigma (\mathbf{F}',\mathbf{F}'')\) .

(Use the fact that every bounded subset of F is relatively compact in \(F''\) for \(\sigma(F'', F')\) . To see that \(\delta\) implies \(\gamma\) ), consider first the case that G is complete, and use Cor. 1 of Th. 2 of TVS, III, §3, No.~6. To pass to the general case, embed G in its completion and observe that every point of \(F''\) is in the closure, for the topology \(\sigma(F'', F')\) , of a bounded subset of F.)

¶ 12) a) Let K be a compact space, E = E(K) the subspace of R \(^{K}\) formed by the linear combinations of characteristic functions of open sets; one identifies E with a subspace of the bidual of C(K). Show that a Cauchy sequence for σ(M(K),E) is convergent for σ(M(K),(M(K))') (make use of Prop. 12 and Exer. 15 of Ch. V, §5). From this, deduce that every subset A of M(K) that is relatively compact for the topology σ(M(K),E), is relatively compact for the topology σ(M(K),(M(K))'). (Reduce to the case that K is metrizable, using Exer. 10 b). Show that A is bounded, by applying Exer. 13 of Ch. V, §5. Note that on A, the vague topology σ(M(K),C(K)) and the topology σ(M(K), \(\overline{E}\) ) are identical, \(\overline{E} \supset C(K)\) being the closure of E in the strong dual of M(K); deduce from this, using the fact that C(K) is separable, that every sequence of points of A has a subsequence that is Cauchy for σ(M(K),E); conclude by invoking Eberlein's theorem.)

\begin{enumerate}
\def\labelenumi{\alph{enumi})}
\setcounter{enumi}{1}
\item
  Let T be a locally compact space, m a vectorial measure on T with values in a quasi-complete, Hausdorff locally convex space F. Assume that, for every compact subset K of T, \(\int_{K} dm \in F\) ; show that, for every bounded Borel f function on T with compact support, one has \(\int f dm \in F\) , and that the image under \(f \mapsto \int f dm\) of the set of Borel functions with support contained in a compact subset K of T and of norm \(\|f\| \leqslant 1\) , is relatively compact for \(\sigma(F, F')\) (make use of a) and Exer. 11, applied to \(u : f \mapsto \int f dm\) , where f runs over the set of linear combinations of characteristic functions of compact subsets of T).
\item
  Let F be a quasi-complete, Hausdorff locally convex space, m a vectorial measure on T with values in F, scalarly of base \(\mu\) ; for every \(z^{\prime} \in F^{\prime}\) , write \(z^{\prime} \circ m = g_{z^{\prime}} \cdot \mu\) . For the condition of b) to be fulfilled, it is necessary and sufficient that, for every compact subset K of T, every equicontinuous subset \(H^{\prime}\) of \(F^{\prime}\) , and every \(\varepsilon > 0\) , there exist a \(\delta > 0\) such that the relations \(A \subset K\) and \(\mu^{*}(A) \leqslant \delta\) imply \(\int_{A}^{*}|g_{z^{\prime}}| d\mu \leqslant \varepsilon\) for every \(z^{\prime} \in H^{\prime}\) (use Exer. 11 above and Exer. 15 of Ch. V, §5); m is then said to be absolutely continuous with respect to \(\mu\) (for the original topology of F). Every vectorial measure with values in a semi-reflexive space and scalarly of base \(\mu\) , is absolutely continuous with respect to \(\mu\) (Cor. 2 of Prop. 3). Every majorizable vectorial measure m with values in a Banach space \(\mathbf{F}\) is absolutely continuous with respect to \(|\mathbf{m}|\) (for the original topology of \(\mathbf{F}\) ).
\end{enumerate}

¶ 13) Let G be a Hausdorff barreled space, F = G' its dual; if, when F is equipped with a topology compatible with the duality between F and G, m is a vectorial measure on T with values in F, then m is again a vectorial measure when F is equipped with the strong topology of the dual of G.

Assume m to be scalarly of base \(\mu\) ; then m is absolutely continuous with respect to \(\mu\) , when F is equipped with the topology \(\tau(\mathrm{F},\mathrm{G})\) (Exer. 12 c)). In order that m be absolutely continuous with respect to \(\mu\) when F is equipped with the strong topology, it is necessary and sufficient that, for every compact subset K of T and every sequence \((\mathrm{A}_{n})\) of \(\mu\) -measurable subsets of K, the image B under m of the closure in \(\mathcal{L}^{1}(\mu)\) of the set of step functions of absolute value \(\leqslant1\) , over the clan generated by the sequence of the \(A_{n}\) (Ch. IV, §4, No.~9), contain a dense countable subset for the strong topology of \(G'\) . (To see that the condition is sufficient, reduce to the case that T = K is a Stone space (Ch. IV, §4, Exer. 10). Then argue by contradiction: if \((\mathrm{A}_{n})\) is a sequence of subsets of K that are both open and closed, one can assume, on passing to a metrizable quotient space \(K_{1}\) of K (Exer. 10 a)), that the continuous functions on \(K_{1}\) corresponding to the \(\varphi_{A_{n}}\) form a total set in \(\mathcal{C}(\mathrm{K}_{1})\) . On the other hand, on considering the quotient of G by the subspace orthogonal to B, one can suppose that B is total in \(F = G'\) for the topology \(\sigma(\mathrm{G}',\mathrm{G})\) ; if H is the subspace of \(F = G'\) generated by B, the hypothesis then implies that every bounded subset C of G is precompact and metrizable for \(\sigma(\mathrm{G},\mathrm{H})\) . Therefore every bounded sequence \((\mathbf{z}_{n})\) of points of G has a subsequence Cauchy for \(\sigma(\mathrm{G},\mathrm{H})\) ; show that if \((\mathbf{z}_{n})\) is a Cauchy sequence for \(\sigma(\mathrm{G},\mathrm{H})\) , and if one writes \(z_{n} \circ m = g_{z_{n}} \cdot \mu\) , the sequence of classes of the \(g_{z_{n}}\) in \(L^{1}(\mu)\) is a Cauchy sequence for the topology \(\sigma(\mathrm{L}^{1},\mathrm{L}^{\infty})\) ; finally, conclude a contradiction from this by using Exers. 15 and 16 of Ch. V, §5.)

\begin{enumerate}
\def\labelenumi{\arabic{enumi})}
\setcounter{enumi}{13}
\tightlist
\item
  \begin{enumerate}
  \def\labelenumii{\alph{enumii})}
  \tightlist
  \item
    Let G be a Banach space, \(F = G'\) its strong dual, g a mapping of T into F belonging to the space \(\Lambda_{G'}^{p}\) (§1, Exer. 16). Show that if p \textgreater{} 1, the measure \(g \cdot \mu\) is absolutely continuous with respect to \(\mu\) (for the strong topology of F).
  \end{enumerate}
\end{enumerate}

\begin{enumerate}
\def\labelenumi{\alph{enumi})}
\setcounter{enumi}{1}
\tightlist
\item
  Take for G the Banach space \(\mathrm{L}^{1}(\mathbf{N})\) , so that \(\mathrm{G}^{\prime}=\mathrm{L}^{\infty}(\mathbf{N})\) ; if f is the function defined in Exer. 7a) of §1 (regarded as taking its values in \(G^{\prime}\) ), then f is scalarly well integrable but the measure \(f\cdot\mu\) is not absolutely continuous with respect to \(\mu\) for the strong topology of \(G^{\prime}\) .
\end{enumerate}

¶ 15) a) Let f be a scalarly locally integrable mapping of T into a quasi-complete Hausdorff locally convex space F. Then \(m = f \cdot \mu\) is a vectorial measure with values in \(F'^{*}\) (for the topology \(\sigma(F'^{*}, F')\) ). Show that if this measure is absolutely continuous (for the topology of uniform convergence on the equicontinuous subsets of \(F'\) ), and if, moreover, f is \(\mu\) -measurable (for the original topology of F), then m is a vectorial measure with values in F. (Make use of Prop. 8 of §1, No.~2.) If F is a Banach space and if, for some p \textgreater{} 1, \(\langle z', f \rangle\) belongs to \(\overline{\mathcal{L}}^{p}(\mu)\) for all \(z' \in F'\) , then the measure m is absolutely continuous (cf.~Exer. 14 a)).

\begin{enumerate}
\def\labelenumi{\alph{enumi})}
\setcounter{enumi}{1}
\tightlist
\item
  Let \(\mathrm{I} = [0,1]\) , \(f\) the function defined on \(\mathrm{I} \times \mathrm{I}\) (using the continuum hypothesis) such that, for every \(t \in \mathrm{I}\) , \(s \mapsto f(s,t)\) is the characteristic function of a countable set, and, for every \(s \in \mathrm{I}\) , \(t \mapsto f(s,t)\) is the characteristic function of the complement of a countable set (Ch. V, §8, Exer. 7 c)). Denote by \(\mathrm{I}_0\) the interval I equipped with the discrete topology, by F the Banach space \(\mathcal{B}(\mathrm{I}) = \mathrm{L}^\infty(\mathrm{I})\) of bounded functions on I; one can identify F with the space \(\mathcal{C}(\widetilde{\mathrm{I}}_0)\) of continuous functions on the Stone-Čech compactification of \(\mathrm{I}_0\) . For every \(s \in \mathrm{I}\) , denote by \(\mathbf{g}(s)\) the element \(t \mapsto f(s,t)\) of F; show that \(\mathbf{g}\) is scalarly integrable for the Lebesgue measure \(\mu\) on I; but \(\int \mathbf{g} d\mu\) does not belong to F. More precisely, \(\mathbf{g} \cdot \mu\) is a vectorial measure with values in \(\mathrm{F}''\) , equal to \(\mathbf{c}\mu\) , where \(\mathbf{c}\) is the constant vector identified with the characteristic function of \(\widetilde{\mathrm{I}}_0 - \mathrm{I}_0 = \mathrm{E}\) (cf.~Exer. 12 a)). (Decompose every positive measure on \(\widetilde{\mathrm{I}}_0\) as the sum of an atomic measure and a diffuse measure (Ch. V, §5, No.~10, Prop. 15), and observe that the diffuse measure is carried by E.) From this, deduce that there does not exist any \(\mu\) -measurable function \(\mathbf{h}\) with values in \(\mathbf{F}\) such that \(\mathbf{h} - \mathbf{g}\) is scalarly \(\mu\) -negligible (make use of \(a\) )).
\end{enumerate}

\begin{enumerate}
\def\labelenumi{\arabic{enumi})}
\setcounter{enumi}{15}
\item
  Consider the mappings \(u_{m}\) of T = {[}0,1{]} into a separable Hilbert space, defined in §1, Exer. 1, and set \(f_{k} = u_{1} + \cdots + u_{k}\) ; show that the vectorial measures \(f_{k} \cdot \mu\) converge uniformly on every bounded subset of \(\mathcal{C}(T)\) to a vectorial measure m with values in F. Moreover, m may be extended by continuity to the space \(\mathcal{L}^{2}(\mu)\) and, for every function \(f \in L^{2}\) , one has \(\int f dm \in F\) and \(|\int f dm| \leqslant N_{2}(f)\) in F; in particular, m is scalarly of base \(\mu\) and absolutely continuous with respect to \(\mu\) for the strong topology (Exer. 12 c)). However, m is not of base \(\nu\) for any positive measure \(\nu\) on T, and a fortiori (No.~5, Cor. 4 of Th. 1) is not majorizable. (Reduce to the case that \(\nu\) has base \(\mu\) , using Th. 3 of Ch. V, §5, No.~7.)
\item
  Let \(\mu\) be Lebesgue measure on \(T = [0,1]\) ; the mapping \(f \mapsto f \cdot \mu\) of \(\mathcal{C}(T)\) into \(\mathcal{M}(T)\) is continuous for the strong topology on \(\mathcal{M}(T)\) , hence is a vectorial measure \(\mathbf{m}\) with values in that Banach space. Show that \(\mathbf{m}\) is scalarly of base \(\mu\) and absolutely continuous with respect to \(\mu\) for the strong topology, but that \(\mathbf{m}\) does not have base \(\mu\) (for the strong topology on \(\mathcal{M}(T)\) ). (Observe that \(\mathbf{m}\) has base \(\mu\) for the vague topology \(\sigma(\mathcal{M}(T), \mathcal{C}(T))\) , and deduce from this that if one had \(\mathbf{m} = \mathbf{g} \cdot \mu\) , with \(\mathbf{g}\) scalarly \(\mu\) -integrable (for the strong topology on \(\mathcal{M}(T)\) ), one would necessarily have \(\mathbf{g}(t) = \varepsilon_t\) almost everywhere. Show that this leads to a contradiction, by observing that for every bounded numerical function \(\theta\) on \(T\) , the linear form \(\mathbf{z}' : \lambda \mapsto \sum_{t \in T} \theta(t) \lambda(\{t\})\)
\end{enumerate}

on \(\mathcal{M}(\mathrm{T})\) is continuous but that \(\langle \mathbf{g},\mathbf{z}'\rangle\) is not necessarily \(\mu\) -measurable.)

¶ 18) a) Let T be a locally compact space, \((K_{\alpha})\) a locally finite covering of T by compact sets. Let \(\mu\) be a positive measure on T, \(\mu_{\alpha}\) the measure induced by \(\mu\) on \(K_{\alpha}\) . Assume that each of the spaces \(L^{\infty}(K_{\alpha}, \mu_{\alpha})\) has the lifting property. Show that \(L^{\infty}(T, \mu)\) has the lifting property.

\begin{enumerate}
\def\labelenumi{\alph{enumi})}
\setcounter{enumi}{1}
\tightlist
\item
  Let K be a metrizable compact space, \(\mu\) a positive measure on K, \((\varpi_n)\) a fundamental sequence (Ch. IV, §5, Exer. 17) of finite partitions of K into integrable sets. For every integrable subset A of K and every element \(\widetilde{f} \in \mathrm{L}^1(\mu)\) , set \(\lambda_A(\widetilde{f}) = 0\) if \(\mu(A) = 0\) , and \(\lambda_A(\widetilde{f}) = (\mu(A))^{-1} \int_{A} f d\mu\) otherwise; for every \(n\) , set \(\rho_n(\widetilde{f}) = \sum_k \lambda_{A_k}(\widetilde{f}) \varphi_{A_k}\) , if
\end{enumerate}

\(\varpi_{n} = (\mathrm{A}_{k})\) . Show that the sequence \((\rho_{n}(\tilde{f}))\) of \(\mu\) -integrable functions converges almost everywhere to \(f\) . (Consider first the case that \(f\) is a continuous function. Given an arbitrary number \(a > 0\) , show next that the union B of all the sets A belonging to at least one of the partitions \(\varpi_{n}\) and such that \(a \cdot \mu(A) \leqslant \int_{A} f d\mu\) , is measurable and is such that \(\mu(B) \leqslant a^{-1} \int f d\mu\) . Then approximate \(f\) in \(\mathcal{L}^{1}\) by a sequence of continuous functions.)

\begin{enumerate}
\def\labelenumi{\alph{enumi})}
\setcounter{enumi}{2}
\item
  Show that every metrizable compact space has the lifting property. (Let \(\mathfrak{U}\) be an ultrafilter on \(\mathbf{N}\) finer than the Fréchet filter; show that \(\rho(\widetilde{f}) = \lim_{\mathfrak{U}} \rho_n(\widetilde{f})\) is a lifting of \(\mathrm{L}^{\infty}\) .)
\item
  Deduce from \(a\) ) and \(c\) ) that every metrizable locally compact space has (for any positive measure whatever) the lifting property.
\end{enumerate}

\begin{enumerate}
\def\labelenumi{\arabic{enumi})}
\setcounter{enumi}{18}
\tightlist
\item
  Let \(\mu\) be a measure on T such that the Banach space \(L^1 (\mu)\) is separable. Show that every continuous linear mapping of \(L^1 (\mu)\) into the strong dual \(F'\) of an arbitrary Banach space \(F\) may be obtained by passage to the quotient starting from a mapping of the form \(g\mapsto \int g\mathbf{f}d\mu\) , where \(\mathbf{f}\in \mathcal{L}_{\mathrm{F}_s'}^\infty\) . (Reduce to the Dunford-Pettis theorem, using Exer. 19 of TVS, IV, §2.)
\end{enumerate}

¶ 20) Let F be a separable Fréchet space, F' its dual, μ a positive measure on T, m a measure on T with values in the weak dual \(F'_s\), scalarly of base μ. For m to have base \(\mu\) , it is necessary and sufficient that the following condition be fulfilled: for every \(\varepsilon > 0\) and every compact subset \(K\) of \(T\) , there exists a compact subset \(K_1 \subset K\) such that \(\mu(K - K_1) \leqslant \varepsilon\) , and such that the image under \(m\) of the set of \(\mu\) -measurable functions \(g\) that are bounded, have support contained in \(K_1\) and satisfy \(\int |g| d\mu \leqslant 1\) , is an equicontinuous subset of \(F'\) . (Recall that \(\int g dm \in F'\) for every \(\mu\) -measurable function \(g\) that is bounded and has compact support; cf.~Cor. 2 of Prop. 3. To show that the condition is necessary, use Prop. 13 of §1, No.~5 and Prop. 5 of §1, No.~2. To see that the condition is sufficient, consider first the case that \(T\) is compact; applying Cor. 3 of Th. 1, No.~5, first show that there exist a partition of \(T\) into a negligible set \(N\) and a sequence ( \(K_n\) ) of compact sets, and a measurable mapping \(g\) of \(T\) into \(F'_s\) , such that \(\int f dm = \int g f d\mu\) for every function \(f \in \mathcal{L}^\infty(\mu)\) that is zero on the complement of the union of a finite number of the sets \(K_n\) ; to prove that \(m = g \cdot \mu\) , use Exers. 16 b) and 20 a) of Ch. V, §5. Finally, to pass to the case that \(T\) is any locally compact space, use Prop. 14 of Ch. IV, §5, No.~9.)

¶ 21) Let F be a Banach space, u a continuous linear form on the Banach space \(\mathrm{L}_{\mathrm{F}}^{p}(\mathrm{T},\mu)\) ; denote by q the conjugate exponent of p, by E the space of numerical step functions over the \(\mu\) -integrable sets. For \(f \in L^{p}\) and \(z \in F\) , one can write \(u(fz) = \langle z, m(f) \rangle\) , where m is a continuous linear mapping of \(L^{p}\) into \(F'\) such that \(|m(f)| \leqslant \|u\| \cdot N_{p}(f)\) .

\begin{enumerate}
\def\labelenumi{\alph{enumi})}
\tightlist
\item
  Let \((\mathrm{A}_i)_{1\leqslant i\leqslant n}\) be a finite sequence of \(\mu\) -integrable sets, pairwise disjoint and non-negligible. Show that
\end{enumerate}

\[
\sum_ {i} \left(| \mathbf {m} (\varphi_ {\mathrm{A} _ {i}}) | ^ {q} / (\mu (\mathrm{A} _ {i})) ^ {q - 1}\right) \leqslant \| u \| ^ {q}.
\]

(For every system \((\mathbf{a}_i)_{1\leqslant i\leqslant n}\) of vectors in F such that \(|\mathbf{a}_i| = 1\) , consider the linear form \((\xi_i)\mapsto u\big(\sum_i\xi_i\mathbf{a}_i\varphi_{\mathrm{A}_i}\big)\) on \(\mathbf{R}^n\) , and apply Th. 4 of Ch. V, §5, No.~8 to a measure with finite support.) From this, deduce that if \(\mathrm{A} = \bigcup_{i}\mathrm{A}_{i}\) , then \(\sum_{i}|\mathbf{m}(\varphi_{\mathrm{A}_i})|\leqslant \| u\| (\mu (\mathrm{A}))^{1 / p}\) (apply Minkowski's inequality).

\begin{enumerate}
\def\labelenumi{\alph{enumi})}
\setcounter{enumi}{1}
\tightlist
\item
  For every positive function \(f \in \mathcal{E}\) , set \(\nu(f) = \sup\left(\sum_{i} |\mathbf{m}(f\varphi_{\mathrm{A}_i})|\right)\) , where \((\mathrm{A}_i)\)
\end{enumerate}

runs over the set of finite sequences of pairwise disjoint \(\mu\) -integrable sets. Show that \(\nu\) is the restriction to \(\mathcal{E}\) of a positive measure with base \(\mu\) (again denoted \(\nu\) ) on T, such that \(|\mathbf{m}(f)| \leqslant \nu(f)\) for every positive function \(f \in \mathcal{E}\) (make use of \(a\) )).

\begin{enumerate}
\def\labelenumi{\alph{enumi})}
\setcounter{enumi}{2}
\tightlist
\item
  Show that if F is separable, there exists a mapping g of T into F', scalarly μ-measurable, such that \(u(\widetilde{\mathbf{f}})=\int\langle\mathbf{f},\mathbf{g}\rangle d\mu\) for every function \(f\in L_{F}^{p}\) , and that \(\|u\|=N_{q}(\mathbf{g})\) . (To establish the last equality, use Prop. 13 of §1, No.~5 and Exer. 13 of §1.) d) Deduce from c) that if F is a reflexive separable Banach space, the space \(L_{F}^{p}\) is reflexive for \(1<p<+\infty\) (apply Prop. 12 of §1, No.~5).
\end{enumerate}

¶ 22) Let F be a Hausdorff locally convex space, S the set of subsets of F that are convex, balanced, and compact for \(\sigma(F, F')\) , and \(S'\) the set of subsets of \(F'\) that are convex, equicontinuous, balanced, and compact for \(\sigma(F', F'')\) .

\begin{enumerate}
\def\labelenumi{\alph{enumi})}
\item
  Suppose that every set in \(\mathfrak{S}\) is precompact for the \(\mathfrak{S}'\) -topology. Show that every continuous linear mapping \(u\) of \(F\) into a quasi-complete, Hausdorff locally convex space \(G\) , that transforms the bounded subsets of \(F\) into subsets relatively compact for \(\sigma(G, G')\) , transforms every set belonging to \(\mathfrak{S}\) into a subset of \(G\) that is relatively compact for the original topology. (Use Exer. 4 of TVS, IV, §1 and Exer. 11 above.) Converse (consider the canonical mapping of \(F\) into its completion for the \(\mathfrak{S}'\) -topology).
\item
  Suppose that F is a metrizable space, or a strict direct limit of metrizable spaces. Show that the condition of a) is satisfied if every Cauchy sequence for \(\sigma(\mathrm{F},\mathrm{F}^{\prime})\) is a Cauchy sequence for the \(S^{\prime}\) -topology (make use of Šmulian's theorem (TVS, IV, §5, Exer. 2 c))).
\item
  Suppose F is infra-barreled (TVS, III, §4, Exer. 7) and denote by \(\mathfrak{S}''\) the set of subsets of \(\mathrm{F}''\) that are convex, equicontinuous, balanced, and compact for \(\sigma(\mathrm{F}''\) , \(\mathrm{F}'''\) ). Show that if every set in \(\mathfrak{S}'\) is precompact for the \(\mathfrak{S}''\) -topology, then every set in \(\mathfrak{S}\) is precompact for the \(\mathfrak{S}'\) -topology. (Use Exer. 4 of TVS, IV, §1, and observe that the original topology of F is induced by the strong topology on \(\mathrm{F}''\) .)
\end{enumerate}

¶ 23) Let T be a locally compact space, \(\mu\) a positive measure on T, and F one of the three Banach spaces \(\overline{\mathcal{K}(T)}\) , \(L^1(\mu)\) , \(L^\infty(\mu)\) . Let u be a continuous linear mapping of F into a quasi-complete space G, that transforms the unit ball of F into a subset of G that is relatively compact for \(\sigma(G, G')\) . Show that u transforms every subset of F relatively compact for \(\sigma(F, F')\) into a subset of G relatively compact for the original topology. (Apply Exer. 22 c) to reduce the case \(F = L^1(\mu)\) to the case \(F = L^\infty(\mu)\) ; use Exer. 13 of Ch. II, §1 to reduce the case \(L^\infty(\mu)\) to the case \(F = \mathcal{C}(S)\) , where S is a suitable compact space. If \(F = \overline{\mathcal{K}(T)}\) , apply Exer. 22 b), then use Exers. 22 b) and 15 of Ch. V, §5; note that a Cauchy sequence for \(\sigma(F, F')\) converges pointwise on T, and a fortiori converges in measure for every bounded measure on T.)

¶ 24) Let \(\mu\) be a positive measure on \(T\) , \(u\) a linear mapping of \(L^1(\mu)\) into a Banach space \(F\) , transforming the unit ball of \(L^1(\mu)\) into a subset of \(F\) relatively compact for \(\sigma(F, F')\) . Show that there exists a \(\mu\) -measurable mapping \(f\) of \(T\) into \(F\) such that \(|\mathbf{f}(t)| \leqslant \|u\|\) for all \(t \in T\) and such that

\[
u (\widetilde {g}) = \int \mathbf {f} g d \mu
\]

for every function \(g \in \overline{\mathcal{L}}^1(\mu)\) (Dunford-Pettis-Phillips theorem). (First reduce to the case that T is compact, with the help of Prop. 14 of Ch. IV, §5, No.~9. The unit ball B of \(L^\infty(\mu) \subset L^1(\mu)\) is then relatively compact for \(\sigma(L^1, L^\infty)\) (Ch. V, §5, Exer. 15); using Exer. 23, show that the closed linear subspace of F generated by \(u(L^1)\) is separable. One can therefore reduce to the case that F is separable and \(\overline{u(L^1)} = F\) . The closure A in F of the image under \(u\) of the unit ball of \(L^1\) is then compact for \(\sigma(F, F')\) (Ch. IV, §7, Exer. 10); show that it is metrizable for this topology (TVS, III, §3, No.~4, Cor. 2 of Prop. 6). Observe that A may then be identified with the unit ball of the dual of the normed space G obtained by equipping \(F'\) with the norm equal to the gauge of \(A^\circ\) . Show that G is separable, and thus reduce to applying the Dunford-Pettis theorem (Cor. 2 of Th. 1).

¶ 25) Let \(\mu\) be a positive measure on T.

\begin{enumerate}
\def\labelenumi{\alph{enumi})}
\item
  Let \(\mathbf{f}\) be a scalarly \(\mu\) -measurable mapping of \(\mathrm{T}\) into a Banach space \(\mathrm{F}\) , such that \(\mathbf{f}(\mathrm{T})\) is relatively compact for \(\sigma(\mathrm{F},\mathrm{F}')\) . Show that there exists a \(\mu\) -measurable mapping \(\mathbf{g}\) of \(\mathrm{T}\) into \(\mathrm{F}\) such that \(\mathbf{f} - \mathbf{g}\) is scalarly locally negligible. (Consider the mapping \(\widetilde{h} \mapsto \int \mathbf{f}h d\mu\) of \(\mathrm{L}^1(\mu)\) into \(\mathrm{F}\) , and apply Exer. 24 to it, also using Exer. 10 of Ch. IV, §7.)
\item
  Let \(\mathbf{f}\) be a scalarly \(\mu\) -measurable mapping of T into a reflexive Fréchet space F. Show that there exists a \(\mu\) -measurable mapping \(\mathbf{g}\) of T into F such that \(\mathbf{f} - \mathbf{g}\) is scalarly locally negligible (cf.~Exer. 15 b)). (Reduce to the case that T is compact, using Prop. 14 of Ch. IV, §5, No.~9. Next, apply Exer. 23 c) of §1 to reduce to the case that \(\mathbf{f}(T)\) is bounded in F. Then embed F in a countable product of Banach spaces, and use a).)
\item
  Let \(\mathbf{f}\) be a mapping of \(\mathrm{T}\) into a Fréchet space \(\mathrm{F}\) , \(\mu\) -measurable for the topology \(\sigma(\mathrm{F}, \mathrm{F}')\) . Show that \(\mathbf{f}\) is \(\mu\) -measurable for the original topology of \(\mathrm{F}\) . (Reduce to the case that \(\mathrm{T}\) is compact and \(\mathbf{f}\) is continuous for the topology \(\sigma(\mathrm{F}, \mathrm{F}')\) ; conclude as in \(b\) .)
\end{enumerate}

¶ 26) Let S, T be two compact spaces, f a finite numerical function defined on S × T.

\begin{enumerate}
\def\labelenumi{\alph{enumi})}
\item
  In order that the mapping \(s \mapsto f(s, \cdot)\) of S into \(\mathbf{R}^{\mathrm{T}}\) be a continuous mapping of S into the space \(\mathcal{C}(\mathrm{T})\) equipped with the topology \(\sigma(\mathcal{C}(\mathrm{T}), \mathcal{M}(\mathrm{T}))\) , it is necessary and sufficient that \(f\) be bounded and that each of the partial mappings \(f(s, \cdot), f(\cdot, t)\) ( \(s \in \mathrm{S}, t \in \mathrm{T}\) ) be continuous. (To see that the condition is sufficient, first show that the image M of S under \(s \mapsto f(s, \cdot)\) is compact for \(\sigma(\mathcal{C}(\mathrm{T}), \mathcal{M}(\mathrm{T}))\) . For this, observe that \(s \mapsto f(s, \cdot)\) is continuous when \(\mathcal{C}(\mathrm{T})\) is equipped with the topology of pointwise convergence on T; use Eberlein's theorem (TVS, IV, §5, No.~3, Th. 1) to reduce to proving that every sequence \((f(s_n, \cdot))\) has a cluster point in \(\mathcal{C}(\mathrm{T})\) for \(\sigma(\mathcal{C}(\mathrm{T}), \mathcal{M}(\mathrm{T}))\) . Reduce to the case that T is metrizable, by considering a suitable quotient space of T (cf.~Exer. 10 a)), and observe that, on a subset of \(\mathcal{C}(\mathrm{T})\) relatively compact for the topology of pointwise convergence on T, this topology is identical to the topology of pointwise convergence on a dense subset of T. Obtain thus a subsequence of \((f(s_n, \cdot))\) that converges pointwise on T, and apply Lebesgue's theorem. Finally, observe that on M, the topology \(\sigma(\mathcal{C}(\mathrm{T}), \mathcal{M}(\mathrm{T}))\) and the topology of pointwise convergence are identical.)
\item
  Assume that each of the partial mappings \(f(s, \cdot), f(\cdot, t)\) is continuous ( \(s \in S\) , \(t \in T\) ). Show that for every positive measure \(\mu\) on \(S\) and every \(\varepsilon > 0\) , there exists a compact subset \(K \subset S\) such that \(\mu(S - K) \leqslant \varepsilon\) and the restriction of \(f\) to \(K \times T\) is continuous. (Reduce to the case that \(f\) is bounded; apply a) and Exer. 25 c).
\item
  With the same hypotheses as in b), show that f is measurable for every measure \(\nu\) on S × T. (Apply b) to the image of \(\nu\) under the projection of S × T onto S.)
\end{enumerate}

\begin{enumerate}
\def\labelenumi{\arabic{enumi})}
\setcounter{enumi}{26}
\tightlist
\item
  Let \(\mathbf{m}\) be a vectorial measure on \(\mathbf{T}\) with values in a Hausdorff locally convex space \(\mathbf{F}\) .
\end{enumerate}

\begin{enumerate}
\def\labelenumi{\alph{enumi})}
\item
  Show that if the support of \(\mathbf{m}\) is finite, then \(\mathbf{m} = \sum_{i=1}^{n} \mathbf{c}_i \varepsilon_{a_i}\) , where the \(\mathbf{c}_i \in \mathbf{F}\) .
\item
  If \(\mathbf{F}\) is a Banach space, and \(\mathbf{m}\) is continuous for the topology of compact convergence, show that \(\mathbf{m}\) has compact support.
\item
  Give an example of a measure with values in \(\mathbf{R}^{\mathbf{N}}\) , with noncompact support, that is continuous for the topology of compact convergence.
\end{enumerate}

\setcounter{exercisegroup}{3}
\edef\CurrentExerciseGroup{\arabic{exercisegroup}}
\section*{\texorpdfstring{\S\CurrentExerciseGroup}{Section \CurrentExerciseGroup}}
\phantomsection
\addcontentsline{toc}{section}{\texorpdfstring{Exercises for \S\CurrentExerciseGroup}{Exercises for Section \CurrentExerciseGroup}}

\begin{enumerate}
\def\labelenumi{\arabic{enumi})}
\item
  Assume that the hypotheses of Th. 1 of No.~1 are satisfied. Let \(h\) be a locally \(\mu\) -integrable function such that \(p\) is \((h \cdot \mu)\) -proper. Show that the function \(b \mapsto g(b) = \int h(t) d\lambda_b(t)\) , defined locally almost everywhere in B (for the measure \(\nu = p(\mu)\) ) is such that \(p(h \cdot \mu) = g \cdot \nu\) .
\item
  Let B be a locally compact space, \((\nu_{\iota})_{\iota\in I}\) the family of all positive measures on B. Let T be the product space \(I\times B\) , I being equipped with the discrete topology, and let \(\nu\) be the measure on T whose restriction to \(\{\iota\}\times B\) is the canonical image of the measure \(\nu_{\iota}\) on B, for all \(\iota\in I\) . Let p be the projection of T onto B; show that if there exists an uncountable compact subset of B, then \(\nu\) does not admit a pseudo-image measure under p.~(Show that every point of B would have measure \textgreater0 for such a measure.)
\item
  Give an example of a positive measure \(\mu\) on a Polish locally compact space \(T\) and a continuous mapping \(p\) of \(T\) into a Polish locally compact space \(B\) , such that \(p\) is not \(\mu\) -proper.
\item
  Let \(\mathrm{T}\) be the interval {[}0,1{]} of \(\mathbf{R}\) , \(\mu\) the Lebesgue measure on \(\mathrm{T}\) . Let \(\mathrm{R}\{x,y\}\) be the equivalence relation \(x - y \in \mathbf{Q}\) in \(\mathrm{T}\) . Show that \(\mathrm{R}\) is not \(\mu\) -measurable (apply Th. 3 of No.~4), but the graph of \(\mathrm{R}\) in \(\mathrm{T} \times \mathrm{T}\) is negligible for the product measure \(\mu \otimes \mu\) .
\item
  Let T be the union of Cantor's triadic set \(K \subset [0,1]\) and the interval I = {[}1,2{]} of R, and let \(\mu\) be the measure induced on the compact space T by Lebesgue measure. Let P be a nonmeasurable subset of I having the power of the continuum (Ch. IV, §4, Exer. 8), and let \(\psi\) be a bijection of K onto P. Consider the equivalence relation R in T for which every point x not belonging to \(K \cup P\) is its own equivalence class, and the class of a point \(y \in K\) consists of y and \(\psi(y)\) . Show that R is \(\mu\) -measurable, but the saturation of K for R is not \(\mu\) -measurable.
\item
  Let T be a Polish locally compact space, \(\mu\) a positive measure on T, R a \(\mu\) -measurable equivalence relation in T, p a \(\mu\) -measurable mapping of T into a Polish locally compact space B such that \(p(x) = p(y)\) is equivalent to \(R\{x, y\}\) , \(\nu\) a pseudo-image measure of \(\mu\) under p, and \(b \mapsto \lambda_b\) a disintegration of \(\mu\) by R. For every \(b \in B\) such that \(\lambda_b \neq 0\) , let \(C(b)\) be the class mod R that carries \(\lambda_b\) ; show that \(\varphi_{C(b)} \cdot \mu\) is proportional to \(\lambda_b\) . Give an example where \(\varphi_{C(b)} \cdot \mu = 0\) for every \(b \in B\) .
\item
  Let T be a Polish locally compact space, \(\mu\) a bounded positive measure on T, and R a \(\mu\) -measurable equivalence relation in T. Let \(\Omega\) be the subset of the space \(\mathcal{M}(\mathrm{T})\) of measures on T, formed by the measures \(\lambda \geqslant 0\) , of total mass \(\leqslant 1\) and nonzero; \(\Omega\) is locally compact when it is equipped with the vague topology (Ch. III, §1, No.~9, Cor. 2 of Prop. 15). Show that there exists one and only one positive measure \(\rho\) on \(\Omega\) such that: \(1^{\circ} \mu = \int_{\Omega} \lambda d\rho(\lambda)\) ; \(2^{\circ} \rho\) is concentrated on a subset \(B_{0}\) of \(\Omega\) whose elements are measures of total mass 1, carried by the classes mod R, two distinct elements of \(B_{0}\) being carried by distinct classes. (Consider a disintegration \(b \mapsto \lambda_{b}\) of \(\mu\) by the relation R, such that all of the \(\lambda_{b}\) have total mass 1, and the image of B under the mapping \(b \mapsto \lambda_{b}\) of B into \(\Omega\) ; make use of Th. 4.)
\item
  Let T be the interval \([0,2]\) in R, \(\mu\) the Lebesgue measure on T. Let A be a non-Borel set contained in the Cantor set \(K \subset [0,1]\) (cf.~GT, IV, §2, No.~4, Example and IX, §6, Exer. 6). Let S be the equivalence relation in T whose equivalence classes are the sets \(\{x\}\) for \(x \notin A \cup (A + 1)\) and the sets \(\{x, x + 1\}\) for \(x \in A\) . Show that S is \(\mu\) -measurable but that there does not exist any Borel section for S.
\item
  Let \(\mathbf{K}\) be a metrizable compact space, \(f\) a continuous mapping of \(\mathbf{K}\) into a Hausdorff topological space \(\mathbf{E}\) , and \(k\) any integer \(> 1\) . Denote by \(\mathrm{B}_k\) the subset of \(\mathbf{E}\) formed by the \(y\) such that \(\stackrel{-1}{f(y)}\) contains at least \(k\) distinct points; show that \(\mathrm{A}_k = \stackrel{-1}{f(\mathrm{B}_k)}\) is a Borel set in \(\mathbf{K}\) (for every integer \(n > 0\) , let \(\mathrm{B}_{kn}\) be the subset of \(\mathbf{E}\) formed by the \(y\) such that \(\stackrel{-1}{f(y)}\) contains at least \(k\) points whose mutual distances are all \(\geqslant 1/n\) ; show that \(\mathrm{B}_{kn}\) is closed).
\end{enumerate}

¶ 10) Let K be a metrizable compact space, μ a positive measure on K, and f a μ-measurable mapping of K into a Hausdorff topological space E. Denote by I (resp. U) the subset of E formed by the y such that \(\stackrel{-1}{f(y)}\) is infinite (resp. uncountable).

\begin{enumerate}
\def\labelenumi{\alph{enumi})}
\item
  Show that there exists in K a \(\mu\) -measurable set H such that \(f(\mathrm{H}) = \mathrm{I}\) and such that, for every \(y \in \mathrm{E}\) , \(\stackrel{-1}{f(y)} \cap \mathbb{C}\mathrm{H}\) is finite. (Let \((\mathrm{K}_n)\) be an increasing sequence of compact subsets of K such that the restriction of \(f\) to each \(\mathrm{K}_n\) is continuous and the complement N of the union F of the \(\mathrm{K}_n\) is \(\mu\) -negligible; for every \(k > 1\) , let \(\mathrm{B}_k\) be the subset of E formed by the \(y\) such that \(\mathrm{F} \cap \stackrel{-1}{f(y)}\) contains at least \(k\) distinct points, and let \(\mathrm{A}_k = \mathrm{F} \cap \stackrel{-1}{f(\mathrm{B}_k)}\) ; take for H the union of the set \(\mathrm{A} = \bigcap_k \mathrm{A}_k\) and the set \(\mathrm{N} \cap \stackrel{-1}{f(\mathrm{I})}\) , and use Exer. 9.)
\item
  Let \(\mathfrak{F}\) be the set of \(\mu\) -measurable sets \(M \subset H\) such that \(M \cap f(y)\) is finite for every \(y \in E\) ; let \(\alpha\) be the supremum of the measures \(\mu(M)\) for \(M \in \mathfrak{F}\) ; there exists an increasing sequence \((\mathrm{M}_n)\) of sets in \(\mathfrak{F}\) such that \(\lim_{n\to \infty}\mu (\mathrm{M}_n) = \alpha\) . Let \(\mathrm{P} = \bigcup_{n}\mathrm{M}_{n}\)
\end{enumerate}

and let \(S\) be a measurable section of \(H \cap \mathbb{C}P\) for the equivalence relation \(f(x) = f(y)\) (Th. 3); show that \(\mu(S) = 0\) .

\begin{enumerate}
\def\labelenumi{\alph{enumi})}
\setcounter{enumi}{2}
\tightlist
\item
  Show that there exists a \(\mu\) -negligible set \(Z \subset K\) such that \(f(Z) = U\) and that, for every \(\varepsilon > 0\) , there exists a \(\mu\) -measurable set \(L \subset K\) such that \(\mu(L) \leqslant \varepsilon\) and \(f(L) = I\) (observe that \(f(H \cap \mathbb{C}M_n) = I\) for every \(n\) and that \(U \subset f(S)\) ).
\end{enumerate}

¶ 11) Let K be a metrizable compact space, μ a positive measure on K, f a μ-measurable mapping of K into a locally compact space T, and ν a positive measure on T; I and U have the same meaning as in Exer. 10.

\begin{enumerate}
\def\labelenumi{\alph{enumi})}
\item
  Suppose that for every \(\mu\) -negligible set \(N \subset K\) , \(f(N)\) is \(\nu\) -negligible. Then \(U\) is \(\nu\) -negligible and, for every \(\mu\) -measurable subset \(A \subset K\) , \(f(A)\) is \(\nu\) -measurable.
\item
  In order that I be \(\nu\) -negligible and the image under \(f\) of every \(\mu\) -negligible set be \(\nu\) -negligible, it is necessary and sufficient that \(f\) satisfy the following condition: for every \(\varepsilon > 0\) , there exists a \(\delta > 0\) such that, for every subset \(A \subset K\) that is \(\mu\) -measurable and is such that \(\mu(A) \leqslant \delta\) , \(f(A)\) is \(\nu\) -measurable and \(\nu(f(A)) \leqslant \varepsilon\) . (To see that the condition is necessary, argue by contradiction, by considering a decreasing sequence \((A_n)\) of \(\mu\) -measurable sets whose intersection is \(\mu\) -negligible, but such that \(\nu(f(A_n)) \geqslant \alpha > 0\) ; then take the intersection of the \(f(A_n)\) and \(\mathbb{C}\mathrm{I}\).)
\end{enumerate}

¶ 12) Let E be the non-metrizable compact space obtained by equipping the interval \([-1,+1]\) of R with the topology T defined in Exer. 13 a) of GT, IX, §2; let \(\varphi\) be the mapping \(x \mapsto |x|\) of E onto I = {[}0,1{]} (equipped with the topology induced by that of R).

\begin{enumerate}
\def\labelenumi{\alph{enumi})}
\item
  Show that \(\varphi\) is continuous and that, denoting by \(\mathfrak{F}\) the set of subsets of \(E\) of the form \(\overline{\varphi}(A)\) , where \(A\) is a Borel subset of \(I\) , the Borel subsets of \(E\) are the subsets \(B\) such that there exists an \(M \in \mathfrak{F}\) for which \(B \cap \mathbb{C}M\) and \(M \cap \mathbb{C}B\) are countable (observe that every open set of \(E\) belongs to \(\mathfrak{F}\) ).
\item
  Show that for a numerical function \(f\) defined on \(\mathbf{E}\) to be continuous (for \(\mathcal{T}\) ), it is necessary and sufficient that it be regulated and that \(f(x-) = f((-x)+)\) for every \(x \in \mathbf{E}\) . One therefore defines a positive measure \(\mu\) on \(\mathbf{E}\) by setting \(\int f d\mu = \int_0^1 f(x) dx\) . Show that the image of \(\mu\) under \(\varphi\) is the Lebesgue measure on \(\mathbf{I}\) , and the \(\mu\) -negligible sets are the subsets negligible for Lebesgue measure on \([-1, +1]\) .
\item
  Deduce from \(a)\) and \(b)\) that there does not exist a \(\mu\) -measurable section for the \(\mu\) -measurable equivalence relation \(\varphi(x) = \varphi(y)\) in \(E\) .
\end{enumerate}

\chapter*{HISTORICAL NOTE}
\phantomsection
\addcontentsline{toc}{chapter}{HISTORICAL NOTE}

(N.B. --- The Roman numerals refer to the bibliography at the end of this note.)

With the development of the `vector calculus' in the course of the 19th century, it was current practice to have to integrate vector-valued functions, but as long as it was only a question of functions with values in finite-dimensional spaces, this operation posed no problem. It is only with Hilbert's spectral theory that one meets operations that lead naturally to a more general concept of integral: this theory in effect leads to associating, with every continuous Hermitian form \(\Phi(x,y)\) on a Hilbert space H, a family \((E(\lambda))_{\lambda\in\mathbf{R}}\) of orthogonal projectors having the property that, for every pair \((x,y)\) of vectors of H, the function \(\lambda\mapsto(E(\lambda)x|y)\) is of bounded variation and \(\Phi(x,y)=\int\lambda d(E(\lambda)x|y)\) ; if one associates with \(\Phi\) the Hermitian operator A such that \(\Phi(x,y)=(Ax|y)\) , it was tempting to write the preceding formula as \(A=\int\lambda dE(\lambda)\) . But it was only from about 1935 on, after the introduction by Bochner of the (`strong') integration of a function with values in a Banach space, that one began to be preoccupied with defining the integral of vector-valued functions (or the integral with respect to a vectorial measure) in such a way as to be able to legitimately write formulas such as the preceding one. This extension was accomplished essentially by Gelfand (III), Dunford and Pettis (IV) and (V); their results are stated for Banach spaces, but extend without difficulty to more general locally convex spaces.

The idea of decomposing a volume into `slices' and reducing an integral over the volume to an integral over each slice, followed by a single integration, has been used in Analysis ever since the beginning of the infinitesimal Calculus (the `Calculus of indivisibles' of Cavalieri being nothing more than a first outline of this principle, which could even be traced back to Archimedes (see the Hist. Note for Book IV, Chs. I--III)). But in the classical applications, the `slices' were always of a very special and very regular nature (most often open subsets of analytic surfaces, depending analytically on a parameter); it could scarcely have been otherwise in the absence of a general theory of integration. The general problem of the disintegration of a measure was posed and solved by von Neumann in 1932, in connection with ergodic theory (I); at about the same time (and independently) Kolmogoroff, while laying down axiomatic foundations for the Theory of Probability, was led to define in a general way the concept of `conditional probability' and to prove its existence, a problem essentially equivalent to that of disintegration of a measure (II).

\chapter*{BIBLIOGRAPHY}
\phantomsection
\addcontentsline{toc}{chapter}{BIBLIOGRAPHY}

\begin{enumerate}
\def\labelenumi{(\Roman{enumi})}
\item
  J. von NEUMANN, Zur Operatorenmethode in der klassischen Mechanik, Ann. of Math., (2), 33 (1932), pp.~587--642.
\item
  A. KOLMOGOROFF, Grundbegriffe der Wahrscheinlichkeitsrechnung, Berlin (Springer), 1933.
\item
  I. GELFAND, Abstrakte Funktionen und lineare Operatoren, Mat. Sborn., (N.S.), 4 (1938), pp.~235--284.
\item
  N. DUNFORD, Uniformity in linear spaces, Trans. Amer. Math. Soc., 44 (1938), pp.~305--356.
\item
  N. DUNFORD AND B. PETTIS, Linear operations on summable functions, Trans. Amer. Math. Soc., 47 (1940), pp.~323--392.
\end{enumerate}

\chapter*{Index of notations}
\phantomsection
\addcontentsline{toc}{chapter}{Index of notations}

Reference numbers indicate, in order, the chapter, section and subsection (or, exceptionally, exercise).

Chapter III :

\(\mathcal{C}(\mathrm{X};\mathrm{E}),\mathcal{C}(\mathrm{X}),\mathcal{C}(\mathrm{X},\mathrm{A};\mathrm{E}),\mathcal{K}(\mathrm{X};\mathrm{E}),\mathcal{K}(\mathrm{X}),\mathcal{K}(\mathrm{X},\mathrm{A};\mathrm{E}),\mathcal{K}(\mathrm{X},\mathrm{A}),\mathcal{K}_{+}(\mathrm{X})\) (X a locally compact space, E a topological vector space): III, 1, 1. Supp(f) (f a function with values in a vector space or in \(\overline{\mathbf{R}}\) ): III, 1, 1. \(\mathcal{C}^b (\mathrm{X};\mathrm{E}),\mathcal{C}^0 (\mathrm{X};\mathrm{E}):III,1,2.\) \(\| \mathbf{f}\|\) (f a function with values in a normed space): III, 1, 2. \(\mu (f),\langle f,\mu \rangle ,\int fd\mu ,\int f\mu ,\int f(x)d\mu (x),\int f(x)\mu (x)\) (f a function in \(\mathcal{H}(\mathrm{X};\mathbf{C})\) \(\mu\) a (complex) measure): III, 1, 3. \(\mathcal{M}(\mathrm{X};\mathbf{C}),\mathcal{M}(\mathrm{X}),\mathcal{M}_{\mathfrak{S}}(\mathrm{X};\mathbf{C}),\mathcal{M}_{\mathfrak{S}}(\mathrm{X}):III,1,3.\) \(\varepsilon_{a}:III,1,3.\) \(g\cdot \mu\) (g a function in \(\mathcal{C}(\mathrm{X};\mathbf{C})\) : III, 1, 4. \(\overline{\mu},\mathcal{R}\mu ,\mathcal{I}\mu :III,1,5.\) \(\mathcal{M}(\mathrm{X};\mathbf{R}),\mathcal{M}(\mathrm{X}),\mathcal{M}_{+}(\mathrm{X}):III,1,5.\) \(\mu \leqslant \nu\) ( \(\mu ,\nu\) real measures): III, 1, 5. \(\mu^{+},\mu^{-},|\mu |\) ( \(\mu\) a real measure): III, 1, 5. \(|\mu |\) ( \(\mu\) a complex measure): III, 1, 6. \(\| \mu \|\) ( \(\mu\) a measure): III, 1, 8. \(\mathcal{M}^1 (\mathrm{X},\mathbf{R}),\mathcal{M}^1 (\mathrm{X}):III,1,8.\) \(\mu |Y\) ( \(\mu\) a measure on X, Y an open subspace of X): III, 2, 1.\\
Supp(μ) ( \(\mu\) a measure): III, 2, 2. \(\langle \mathbf{f},\mathbf{z}'\rangle :III,3,1.\) \(\widetilde{\mathcal{K}} (\mathrm{X};\mathrm{E}):III,3,1.\) \(\int \mathbf{f}d\mu ,\int \mathbf{f}\mu ,\int \mathbf{f}(x)d\mu (x),\int \mathbf{f}(x)\mu (x)\) ( \(f\) a function in \(\widetilde{\mathcal{K}} (\mathrm{X};\mathrm{E}))\) : III, 3, 1. \(\int d\mu (y)\int f(x,y)d\lambda (x):III,4,1.\) \(\iint f d\lambda d\mu ,\iint f d\mu d\lambda ,\iint f\lambda \mu ,\iint f\mu \lambda ,\iint f(x,y)d\lambda (x)d\mu (y),\) \(\iint f(x,y)d\mu (y)d\lambda (x),\iint f(x,y)\lambda (x)\mu (y),\iint f(x,y)\mu (y)\lambda (x):III,4,1.\)

\(\lambda \otimes \mu (\lambda ,\mu \text{measures}):\mathrm{III},4,2.\) \(\mu_1\otimes \mu_2\otimes \dots \otimes \mu_n,\bigotimes_{i = 1}^n\mu_i:\mathrm{III},4,4.\) \(\int f d\mu_1d\mu_2\dots d\mu_n,\iint \dots \int f d\mu_1d\mu_2\dots d\mu_n,\int f(\mu_1\otimes \mu_2\otimes \dots \otimes \mu_n),\) \(\iint \dots \int f(x_1,x_2,\ldots ,x_n)d\mu_1(x_1)d\mu_2(x_2)\dots d\mu_n(x_n),\) \(\iint \dots \int f(x_1,x_2,\ldots ,x_n)\mu_1(x_1)\mu_2(x_2)\dots \mu_n(x_n):\mathrm{III},4,4.\) \(\bigotimes_{\lambda \in L}\mu_{\lambda}:\mathrm{III},4,6.\)

\(\varphi_{\mathrm{A}}:\mathrm{IV},1,1.\) \(\mathcal{K}_{+},\mathcal{I}_{+}(\mathrm{X}),\mathcal{I}_{+}:\mathrm{IV},1,1.\) \(\mu^{*}(f)\) (\(\mu\) a positive measure): IV, 1, 1, IV, 1, 3 and IV, 4, Exer. 5.

\(\int^{*}fd\mu,\int^{*}f\mu,\int^{*}f(x)d\mu(x),\int^{*}f(x)\mu(x)\) ( \(f\) a function \(\geqslant 0\), \(\mu\) a positive measure): IV, 1, 3.

\(\mu^{*}(\mathrm{A})\) (A a subset of X, \(\mu\) a positive measure): IV, 1, 2 and IV, 1, 4.

\(\widetilde{\mathbf{f}}:\mathrm{IV},2,4\) and IV, 2, 5.

\(\varphi((\widetilde{\mathbf{f}}_n))\), \(\widetilde{\mathbf{f}} +\widetilde{\mathbf{g}}\), \(\alpha \widetilde{\mathbf{f}},\widetilde{g}\widetilde{\mathbf{f}}:\mathrm{IV},2,4.\) \(\widetilde{f}\leqslant \widetilde{g},\sup_{n}\widetilde{f}_n,\inf_n\widetilde{f}_n,\limsup_{n\to \infty}\widetilde{f}_n,\liminf_{n\to \infty}\widetilde{f}_n\) (\(f,g,f_n\) numerical functions): IV, 2, 6.

\(|\mathbf{z}|\) (z a point of a normed space): IV, 3, 2.

\(|\mathbf{f}|\) (f a function with values in a normed space): IV, 3, 2.

\(\mathrm{N}_p(\mathbf{f},\mu),\mathrm{N}_p(\mathbf{f}),\mathrm{N}_p(\widetilde{\mathbf{f}})\) (\(1\leqslant p<+\infty\)): IV, 3, 2.

\(\mathcal{F}_{\mathrm{F}}(\mathrm{X}),\mathcal{F}_{\mathrm{F}}:\mathrm{IV},3,3.\) \(\mathcal{F}_{\mathrm{F}}^{p}(\mathrm{X},\mu),\mathcal{F}_{\mathrm{F}}^{p}(\mu),\mathcal{F}_{\mathrm{F}}^{p}:\mathrm{IV},3,3.\) \(\mathcal{N}_{\mathrm{F}}:\mathrm{IV},3,3.\) \(\mathcal{K}_{\mathrm{F}},\mathcal{L}_{\mathrm{F}}^{p}(\mathrm{X},\mu),\mathcal{L}_{\mathrm{F}}^{p}(\mu),\mathcal{L}_{\mathrm{F}}^{p},\mathrm{L}_{\mathrm{F}}^{p}(\mathrm{X},\mu),\mathrm{L}_{\mathrm{F}}^{p}(\mu),\mathrm{L}_{\mathrm{F}}^{p},\mathcal{L}^{p},\mathrm{L}^{p}\) (\(1\leqslant p<+\infty\)): IV, 3, 4.

\(\| \widetilde{\mathbf{f}} \| _p\) (\(1\leqslant p < + \infty\)): IV, 3, 4.

\(\mu (\mathbf{f}),\int \mathbf{f}d\mu ,\int \mathbf{f}(x)d\mu (x),\int \mathbf{f}\mu ,\int \mathbf{f}(x)\mu (x),\mu (\widetilde{\mathbf{f}})\) (f a \(\mu\) -integrable function with values in a Banach space): IV, 4, 1.

\(\mu (\mathrm{A})\) (A a \(\mu\) -integrable set): IV, 4, 5.

\(\mathcal{C}'(\mathrm{X};\mathbf{C}):\mathrm{IV},4,8.\) \(\mathcal{E}(\Phi),\mathcal{E}_{\mathrm{F}}(\Phi)\) (\(\Phi\) a clan of sets): IV, 4, 9.

\(\mu_{*}(f)\) (\(f\) a function): IV, 4, Exer. 5.

\[
\mu_ {*} (\mathrm{A}) (\text { A a set }): \mathrm{IV}, 4, \text { Exer. } 7.
\]

\[
\int_ {\mathrm{A}} \textbf {f} d \mu , \int_ {\mathrm{A}} \textbf {f} \mu , \int_ {\mathrm{A}} ^ {*} f d \mu , \int_ {\mathrm{A}} ^ {*} f \mu : \mathrm{IV}, 5, 6.
\]

\[
\mathcal {S} (\mathrm{A}, \mu ; \mathrm{F}), \mathcal {S} _ {\mathrm{F}} (\mathrm{A}, \mu), \mathcal {S} _ {\mathrm{F}} (\mu), \mathcal {S} _ {\mathrm{F}}: \mathrm{IV}, 5, 11.
\]

\[
\boldsymbol {W} (\mathrm{V}, \mathrm{B}, \delta): \mathrm{IV}, 5, 11.
\]

\[
\mathrm{S} (\mathrm{A}, \mu ; \mathrm{F}), \mathrm{S} _ {\mathrm{F}} (\mathrm{A}, \mu), \mathrm{S} _ {\mathrm{F}} (\mu), \mathrm{S} _ {\mathrm{F}}: \mathrm{IV}, 5, 11.
\]

\[
\mathrm{M} _ {\infty} (f), m _ {\infty} (f): \mathrm{IV}, 6, 2.
\]

\[
\mathrm{N} _ {\infty} (\mathbf {f}): \mathrm{IV}, 6, 3.
\]

\[
\mathcal {L} _ {\mathrm{F}} ^ {\infty} (\mathrm{X}, \mu), \mathcal {L} _ {\mathrm{F}} ^ {\infty} (\mu), \mathcal {L} _ {\mathrm{F}} ^ {\infty}, \mathcal {N} _ {\mathrm{F}} ^ {\infty}: \mathrm{IV}, 6, 3.
\]

\[
\| \dot {\mathbf {f}} \| _ {\infty}, \mathrm{N} _ {\infty} (\dot {\mathbf {f}}): \mathrm{IV}, 6, 3.
\]

\[
\mathrm{L} _ {\mathrm{F}} ^ {\infty} (\mathrm{X}, \mu), \mathrm{L} _ {\mathrm{F}} ^ {\infty} (\mu), \mathrm{L} _ {\mathrm{F}} ^ {\infty}, \mathcal {L} ^ {\infty}, \mathrm{L} ^ {\infty}: \mathrm{IV}, 6, 3.
\]

\[
\mathbf {b} _ {\mu}: \mathrm{IV}, 7, 1.
\]

\[
\mathrm{Ch} _ {\mathcal {H}} (\mathrm{X}), \mathrm{Ch} (\mathrm{X}), \check {\mathrm{S}} _ {\mathcal {H}} (\mathrm{X}), \check {\mathrm{S}} (\mathrm{X}): \mathrm{IV}, 7, 3.
\]

\[
\mu^ {\bullet} (f), \mu^ {\bullet} (\mathrm{A}), \int^ {\bullet} f d \mu , \int^ {\bullet} f (t) d \mu (t), \int^ {\bullet} f \mu : \mathrm{V}, 1, 1.
\]

\[
\overline {{\mathcal {F}}} _ {\mathrm{F}} ^ {p} (\mathrm{T}, \mu), \overline {{\mathcal {F}}} _ {\mathrm{F}} ^ {p} (\mu), \overline {{\mathcal {F}}} _ {\mathrm{F}} ^ {p}: \mathrm{V}, 1, 3.
\]

\[
\overline {{\mathbf {N}}} _ {p} (f), \overline {{\mathcal {L}}} _ {\mathrm{F}} ^ {p} (\mathrm{T}, \mu), \overline {{\mathcal {L}}} _ {\mathrm{F}} ^ {p} (\mu), \overline {{\mathcal {L}}} _ {\mathrm{F}} ^ {p}: \mathrm{V}, 1, 3.
\]

\[
\overline {{\mathcal {L}}} _ {\mathrm{F}} ^ {p} (\mathrm{T}, \theta) (\theta \text {a complex measure}): \mathrm{V,1,3.}
\]

\[
\int \lambda_ {t} d \mu (t) (t \mapsto \lambda_ {t} \text {a family of positive measures}): \mathrm{V}, 3, 1.
\]

\[
\int d \mu (t) \int f (x) d \lambda_ {t} (x): \mathrm{V}, 3, 1.
\]

\[
\| \Lambda \| (\Lambda \text {a diffusion}): \mathrm{V}, 3, 5.
\]

\[
\langle \eta , h \rangle : \mathrm{V}, 3, 5.
\]

\[
\Lambda f, \mu \Lambda : \mathrm{V}, 3, 5.
\]

\[
\Lambda \mathrm{H}: \mathrm{V}, 3, 6.
\]

\[
\mathcal {L} _ {\mathrm{loc}} ^ {1} (\mathrm{T}, \mu ; \mathrm{F}), \mathrm{L} _ {\mathrm{loc}} ^ {1} (\mathrm{T}, \mu ; \mathrm{F}): \mathrm{V}, 5, 1.
\]

\[
\int_ {\mathrm{A}} ^ {\bullet} f d \mu : \mathrm{V}, 5, 3.
\]

\[
u (\mu_ {1}, \dots , \mu_ {n}) (u \text {a positively homogeneous numerical function}): \mathrm{V}, 5, 9.
\]

\[
\pi^ {- 1} (\mu) (\pi \text {a local homeomorphism}) \colon \mathrm{V,6,6.}
\]

\[
\iint^ {*} f (t, t ^ {\prime}) d \mu (t) d \mu^ {\prime} (t ^ {\prime}), \iint^ {\bullet} f (t, t ^ {\prime}) d \mu (t) d \mu^ {\prime} (t ^ {\prime}), \iint \mathbf {f} (t, t ^ {\prime}) d \mu (t) d \mu^ {\prime} (t ^ {\prime}): \mathrm{V}, 8, 1
\]

Chapter VI :

\(\mathbf{F}^{\prime},\mathbf{F}^{\prime \prime},\mathbf{F}^{\prime *},\mathbf{F}_{\sigma}\) (F a Hausdorff locally convex space): VI, Introduction.

\(\mathcal{K}(\mathrm{T}),\mathcal{K}_{\mathbf{R}}(\mathrm{T}),\mathcal{K}_{\mathbf{C}}(\mathrm{T}),\mathcal{K}(\mathrm{T},\mathrm{A}),\mathcal{K}_{\mathbf{C}}(\mathrm{T},\mathrm{A}):\mathrm{VI},\) Introduction.

\(\langle \mathbf{f},\mathbf{z}^{\prime}\rangle ,\langle \mathbf{z}^{\prime},\mathbf{f}\rangle :\mathrm{VI},1.\)

\(\int \mathbf{f} d\mu, \int \mathbf{f}(t) d\mu(t)\) (f a vector-valued function, \(\mu\) a positive measure): VI, 1, 1.

\(g\mathbf{f},\mathbf{f}g\) (f a vector-valued function, \(g\) a scalar function): VI, 1, 1.

\(\mathcal{C}'(\mathrm{T}) : \mathrm{VI}, 1, 6.\)

\(\int f d\mathbf{m}, \int f(t) d\mathbf{m}(t)\) ( \(f\) a numerical function, \(\mathbf{m}\) a vectorial measure): VI, 2, 1 and VI, 2, 2.

\(g \cdot m\) (g a numerical function, m a vectorial measure): VI, 2, 1.

\(\mathcal{L}(\mathbf{m}) : \mathrm{VI}, 2, 2.\)

\(q(\mathbf{m}), |\mathbf{m}|\) ( \(q\) a semi-norm, \(\mathbf{m}\) a vectorial measure): VI, 2, 3.

\(\mathbf{f} \cdot \mu\) (f a vector-valued function, \(\mu\) a positive measure): VI, 2, 4.

\(\mathcal{L}_{\mathrm{F}_s'}^{\infty}, \mathrm{L}_{\mathrm{F}_s'}^{\infty} : \mathrm{VI}, 2, 5.\)

\(\langle \mathbf{f},\mathbf{g}\rangle\) (f,g vector-valued functions): VI, 2, 6.

\(\mathrm{I}_{\Phi ,\mathbf{m}},\int \mathbf{f}dm\) (f a vector-valued function, m a vectorial measure): VI, 2, 7.

\(|m|, \int \mathbf{f} dm \ (m \ a \ complex \ measure) : VI, 2, 8, III, 1, 6.\)

\(\mathcal{L}_{\mathrm{F}}^{p}(\mathrm{T},m), \overline{\mathcal{L}}_{\mathrm{F}}^{p}(\mathrm{T},m), \mathrm{L}_{\mathrm{F}}^{p}(\mathrm{T},m) (m \text{ a complex measure}): \mathrm{VI}, 2, 8, \mathrm{V}, 1, 3.\)

\(h \cdot m\) ( \(m\) a complex measure): VI, 2, 8.

\(\overline{m}\) ( \(m\) a complex measure): VI, 2, 8, III, 1, 5.

\(\| m\|\) ( \(m\) a complex measure): VI, 2, 9, III, 1, 8.

\(\pi (m), m_{\mathrm{Y}}, m \otimes m' (m, m'\) complex measures): VI, 2, 10.

\(\mathfrak{B}(\mathrm{F}_{1},\mathrm{F}_{2}),{}^{r}\Phi,{}^{l}\Phi :\mathrm{VI},\mathrm{App.},1.\)

\(\mathrm{E}_{\sigma}, \mathrm{F}_{\sigma}, \mathrm{E}_{s}^{\prime}, \mathrm{F}_{s}^{\prime}, \mathcal{B}(\mathrm{E}, \mathrm{F}) : \mathrm{VI}, \mathrm{App.}, 1.\)

\(\Lambda_{\mathrm{F}^{\prime}}^{p}(\mathrm{T},\mu),\mathrm{M}_{p},\mathrm{M}_{p}^{\prime}:\mathrm{VI},1,\mathrm{Exer.}16.\)
